\documentclass[final,3p,authoryear]{elsarticle}

% Packages required by this manuscript.
\usepackage{amsmath}
\usepackage{amsthm}
\usepackage{amssymb}
\usepackage{amsfonts}
\usepackage{mathrsfs}
\usepackage{mathtools}
\usepackage{graphicx}
\usepackage{enumerate}
\usepackage{enumitem}
\usepackage[table]{xcolor}
\usepackage{natbib}
\usepackage{url}
\usepackage{booktabs}
\usepackage{multirow}
\usepackage{algorithm}
\usepackage{algpseudocode}
\usepackage{hyperref}
\usepackage{tabularx}
\usepackage{longtable}
\usepackage{makecell}
\usepackage{subcaption}
\usepackage[export]{adjustbox}
\usepackage{float}
\usepackage{array}
\usepackage{caption}
\usepackage{threeparttable}
\usepackage{tikz}

\allowdisplaybreaks[1]
\newtheorem{definition}{Definition}
\newtheorem{proposition}{Proposition}

\newlength{\thumbw}
\newlength{\thumbh}
\setlength{\thumbw}{0.225\textwidth}
\setlength{\thumbh}{2.25cm}

\newcommand{\cropthumb}[2]{%
\begin{tikzpicture}
\clip (0,0) rectangle (\thumbw,\thumbh);
\node[anchor=center] at (0.5\thumbw,0.5\thumbh)
{\includegraphics[#1]{#2}};
\end{tikzpicture}%
}


\begin{document}

\begin{frontmatter}

\title{Strategic Design of Urban Drone Delivery Networks: Integrating Multi-Type Facilities, Customer Choice, Congestion, and Fleet Circulation}
\author{qimy}
\date{}

\begin{abstract}
% Previous wording retained for revision history: Urban low-altitude logistics can supplement ground delivery, but its viability depends on a network design that coordinates heterogeneous infrastructure, customer adoption, congestion, energy feasibility, and drone availability.
\color{blue}
Routine deployments have demonstrated that urban low-altitude logistics can be operationally viable as a supplement to ground delivery in selected settings, while its commercial viability at scale remains unresolved.
We therefore study how the joint design of heterogeneous infrastructure under a general mixed-delivery formulation and its home-delivery-assistance and locker-only restrictions, together with customer adoption, congestion, energy feasibility, and drone availability, shapes platform fulfillment cost and service performance.
\color{black}
We study a strategic planning problem in which drones selectively offload courier demand through physical drone bases, \textcolor{blue}{customer-facing smart lockers}, battery-swapping stations, and virtual idle-drone states.
The user-choice module determines the locker-pickup classification of successful service and, in the locker-only restriction, the demand eligible for the drone network, while a closed queueing network represents drone circulation, \textcolor{blue}{smart-locker congestion}, task chaining, and fleet requirements.
We formulate a mixed-integer second-order conic program that uses discrete availability levels and continuous demand-allocation variables to preserve a linear interface between user choice and operations.
% Previous wording retained for revision history: The framework is designed to evaluate the economic and service implications of alternative urban drone-network configurations using operational data from Shenzhen, China.
\color{blue}
Using an anonymized seven-day sample of 2023 delivery orders from Nanshan District, Shenzhen, the framework constructs spatial demand and candidate locations, while customer-choice, cost, and operating parameters remain transparent scenario inputs rather than empirical estimates.
These comparisons inform cost-related conditions relevant to commercial viability but do not establish profitability at scale.
\color{black}
% Urban low-altitude logistics has grown rapidly in recent years and become an important supplement to ground transportation.
% Reasonably deploying drone sites and forming an efficient urban-scale network are key to the sustainable development of low-altitude logistics.
% However, existing studies generally do not distinguish site sub-functions and offer limited flexibility in site configuration.
% To address these challenges, this research introduces a multi-type facility location framework for user-to-user demand scenarios, such as takeaway and same-city express delivery.
% We model customer-oriented smart lockers as unified pickup and delivery interfaces, physical bases as drone parking facilities, and battery-swapping stations as en-route energy support.
% Building on a closed queueing-network approximation, we propose a mixed-integer second-order cone programming model to jointly optimize the location, quantity, capacity, and connectivity of each station type.
% This study contributes to strategic-level optimization for low-altitude logistics network planning and improves the forward-looking and global optimality of drone facility deployment.
\end{abstract}

\begin{keyword}
Network design; Drone-based delivery; Hierarchical facility location; User choice; Closed queueing network
\end{keyword}

\end{frontmatter}

\section{Introduction}
\label{sec:introduction}

% Introduction: Drone Industry

The rapid expansion of e-commerce and the corresponding surge in consumer demand for faster delivery times have placed immense pressure on last-mile logistics.
Traditional ground-based delivery networks are increasingly constrained by urban traffic congestion, leading to higher operational costs and longer delivery windows.
In response, Unmanned Aerial Vehicles (UAVs), commonly known as drones, have emerged as a transformative technology with the potential to revolutionize last-mile delivery.
% Previous unsupported market-size forecast retained for revision history: The global drone logistics and transportation market is projected to experience explosive growth, with some forecasts predicting a compound annual growth rate (CAGR) of over 50\% to reach a market size of over USD 82 billion by 2034.
\textcolor{blue}{Commercial trials and a growing interdisciplinary literature have intensified interest in drone-based last-mile delivery \citep{DroneEconomicViabilityReview2025}.}
% Previous wording retained for revision history: By operating in low-altitude airspace, drones can bypass ground congestion, offering significant reductions in delivery time, operational costs, and carbon emissions.
\color{blue}
By operating in low-altitude airspace, drones can bypass ground congestion and may shorten delivery times in suitable settings.
% Citation link for \citep{DroneEconomicViabilityReview2025}:
% https://doi.org/10.1108/IJLM-04-2023-0149
Whether bypassing congestion translates into lower platform fulfillment cost depends on network design, utilization, demand density, payload, and operating conditions \citep{DroneEconomicViabilityReview2025}.
% Citation links for \citep{GoodchildToy2018DroneEmissions,Stolaroff2018DroneEnergy}:
% https://doi.org/10.1016/j.trd.2017.02.017
% https://doi.org/10.1038/s41467-017-02411-5
Environmental benefits likewise depend on the energy source, payload, vehicle type, and service setting \citep{GoodchildToy2018DroneEmissions,Stolaroff2018DroneEnergy}.
\color{black}

% Introduction: Drone Delivery

This technological shift is particularly salient in the on-demand food delivery sector, a key segment of the broader Urban Air Mobility (UAM) vision.
Companies are actively developing and deploying drone-based solutions to deliver meals quickly while preserving food quality.
% Previous wording retained for revision history: A prominent example is Meituan, one of China's largest on-demand service platforms, which has been operating drone delivery services in dense urban centers like Shenzhen since 2021.
% Citation link for \citep{Meituan2023DroneOperationsHistory}:
% https://www.meituan.com/news/NN230117001043072
\color{blue}
A prominent example is Meituan, one of China's largest on-demand service platforms, which completed its first drone-delivery order in early 2021 and subsequently operated routine trial services in Shenzhen and Shanghai \citep{Meituan2023DroneOperationsHistory}.
% Previous wording retained for revision history: Having completed hundreds of thousands of orders with an average delivery time of approximately 15 minutes, Meituan's operations demonstrate the viability and efficiency of a well-designed drone logistics network.
% Previous wording retained for revision history: This network relies on a sophisticated infrastructure of automated drone bases and customer-facing smart lockers to facilitate its service.
% Citation link for \citep{SFHolding2024FengyiOperations}:
% https://static.cninfo.com.cn/finalpage/2024-04-30/1219968909.PDF
SF Holding reported in 2024 that Fengyi had achieved routine high-volume drone operations in the Greater Bay Area \citep{SFHolding2024FengyiOperations}.
% Citation link for \citep{Meituan2026DroneOperations}:
% https://www.meituan.com/news/NN260522185005269
Meituan reported in 2026 that its low-altitude network had become a stable supplementary delivery network with more than 900,000 completed orders \citep{Meituan2026DroneOperations}.
% Citation link for \citep{FAA2026DroneDeliveryPart135}:
% https://www.faa.gov/uas/advanced_operations/package_delivery_drone
Internationally, Wing, UPS Flight Forward (operating Matternet aircraft), and Zipline received U.S. air-carrier authority under Part 135 and commenced real-world package-delivery operations \citep{FAA2026DroneDeliveryPart135}.
These cases show that drone delivery has moved beyond isolated technical demonstrations into reported routine deployments in China and, for the U.S. operators, Part 135-authorized real-world operations.
Such operating milestones do not, however, establish commercial viability.
% Citation link for \citep{Matternet2026S1}:
% https://www.sec.gov/Archives/edgar/data/2075109/000121390026073335/0001213900-26-073335-index.htm
Matternet reported in its June 2026 S-1 that its platform had supported routine revenue-generating operations but that the company had not been profitable to date \citep{Matternet2026S1}.
% Citation link for \citep{Meituan2026ProfitabilityInterview}:
% https://finance.sina.com.cn/tech/shenji/2026-05-21/doc-inhyrmms9217368.shtml
In a May 2026 interview, the head of Meituan's drone business stated that sustainable profitability at scale remained a future objective \citep{Meituan2026ProfitabilityInterview}.
Here, operational viability denotes the ability to support reliable routine service, whereas commercial viability denotes sustainable profitability at scale.
Across these heterogeneous service architectures, the configuration of bases, customer handoff facilities, energy support, and courier coordination can materially affect cost performance \citep{DroneEconomicViabilityReview2025}.
Because the present model represents cost and service outcomes but not revenue or firm-level profit, it evaluates how network design affects the cost-related conditions that may support commercial viability, without claiming that commercial viability has been achieved.
\color{black}

% Description: Commercial company

% Previous incomplete wording retained for revision history: For sucessful realization of such drone delivery services by commercial companies like Zipline, Wing, Meituan, and Amazon Prime Air,
% 简要介绍，引用数据
% Previous incomplete wording retained for revision history: The

% Description: Problem

% The success of such large-scale operations hinges on the strategic design of the underlying infrastructure.
% This presents a complex strategic decision problem: how to locate and configure a multi-type network that serves distributed user-to-user (U2U) demand while recognizing that customers may choose whether to use the drone-enabled service \citep{contrerasGeneralNetworkDesign2012}.
% % 解释(U2U) 
% The physical network comprises drone bases, unified customer-facing smart lockers that support both pickup and drop-off, and battery-swapping stations that extend operational range.
% A critical operational feature is task chaining, under which a drone that completes a delivery may remain available near a smart locker for a subsequent assignment rather than immediately returning to a physical base.
% Virtual idle-drone states represent this operational availability without being treated as additional physical facilities.
% The resulting system is flow dependent because facility selection affects user choice, realized demand affects drone circulation, and circulation determines congestion and fleet requirements.



\begin{figure}[htbp]
\centering

\begin{tabular}{@{}c@{\hspace{0.015\textwidth}}c@{\hspace{0.045\textwidth}}c@{\hspace{0.015\textwidth}}c@{}}
\cropthumb{height=\thumbh}{Figure/DroneBases.jpg}
&
\cropthumb{height=\thumbh}{Figure/DroneBases2.png}
&
\cropthumb{height=\thumbh}{Figure/Hubs.jpg}
&
\cropthumb{width=\thumbw}{Figure/Hubs2.png}
\\[-0.1em]
\multicolumn{2}{c}{(a) Drone bases}
&
\multicolumn{2}{c}{(b) \textcolor{blue}{Drone lockers}}
\end{tabular}

\caption{Representative infrastructure for urban drone logistics.}
\label{fig:drone_infrastructure}
\vspace{-0.1in}
\end{figure}

% Brief literature review
\color{blue}
Demand-responsive service-network studies illustrate the interdependence of infrastructure, operations, and endogenous demand.
In urban electric-vehicle sharing, \citet{heServiceRegion2017M&SOM} integrate service-region design with customer adoption, fleet repositioning, charging, and uncertainty in customer-adoption patterns.
For urban aerial mobility, Wang et al.~\citeyearpar{wangVertiportPlanning2022M&SOM} jointly optimize vertiport deployment, passenger--vertiport matching, queueing operations, and origin--destination-level adoption.
Pair assignment remains centralized: the operator selects the departure and arrival vertiports and arranges both ground legs, while customer demand remains indexed only by the origin--destination pair and uses the worst-case expected trip time across active pair assignments.

Related work in urban logistics has begun to integrate infrastructure planning with the coordinated operation of drones and couriers.
\citet{chenCoordinatedDroneCourier2025} study a closed drone network among vertiports in which couriers perform the customer-side first- and last-mile legs, and connect vertiport-network planning to queueing-network operations.
This work shows how coordination and asset circulation can be embedded in network planning, but its door-to-door setting takes origin--destination demand and the expected fast-service share as inputs rather than modeling customer choice between home delivery and locker pickup.

These streams leave a specific planning interface to be addressed in our setting.
The destination smart locker $j$ is visible to the customer and affects the home-versus-locker choice, whereas the origin smart locker $i$ is an internal platform assignment that should not create a separate customer-choice alternative.
Infrastructure decisions must therefore determine which eligible destination locker, if any, is displayed, translate its service attributes into a conditional locker-choice share, and coordinate realized drone service with courier fallback, congestion, energy-feasible routing, and closed-fleet circulation.
This distinction is central to the formulation developed below because it connects customer adoption to fulfillment-cost and service consequences without treating backend routing as a customer decision.

% Brief highlights of modeling and solution methods

Guided by these requirements, we aim to address the following critical question:
% Previous wording retained for revision history: \textit{How can an urban drone delivery network and its operating strategies be jointly planned to create business value in the urban on-demand delivery industry?}
\textit{How should an urban drone delivery network be designed to minimize modeled platform fulfillment cost, and how do the resulting cost--service trade-offs inform the cost-related conditions that may support commercial viability in urban on-demand delivery?}
% Previous wording retained for revision history: To answer this question, we develop a unified optimization framework that jointly plans the multi-type infrastructure and incorporates customer preferences for smart lockers.
To investigate this question, we develop a unified optimization framework that jointly plans multi-type infrastructure and service allocation while incorporating customer preferences for smart lockers.
Under a jointly planned queueing network comprising drone bases, smart lockers, and battery-swapping stations, we study a general problem and two restricted service scopes.
(i) The general mixed-delivery problem allows drones to serve both locker-pickup orders and platform-selected home-delivery orders.
(ii) The home-delivery-assistance restriction suppresses the customer-facing locker-pickup option and allows drones to assist only home-delivery orders.
(iii) The locker-only restriction allows customer locker choice but excludes home-delivery orders from drone service.
The latter two restrictions are obtained from the general problem by replacing explicitly identified demand, service, and cost equations and omitting inactive variable and endogenous-quantity families rather than by activating exogenous mode-indicator parameters.
% Any demand not served by drones falls back to couriers.


% Guided by these requirements, we develop an optimization framework for the MLP-IC that links geometric approximation, compact modeling, and decomposition-based exact optimization.
% Specifically, we ask three methodological questions: (i) how to represent irregular continuous-placement geometry with increasing fidelity while preserving a common optimization structure; (ii) how to formulate the resulting loading problem to jointly capture automobile selection, wagon assignment, sequencing, and deck-end configuration decisions; and (iii) how to solve the formulation at scale through deck-consistent decomposition and structure-exploiting pricing.
% Our contributions are as follows.

% Given these considerations, our research question emerges as a critical inquiry: 
% can we model how different locker network designs affect delivery efficiency after incorporating customer preferences in locker usage? 
% Addressing this question requires developing a model that not only accounts for both the fixed locations of lockers but also adapts to the variability in home delivery locations and customer preferences in locker usage. 
% More formally, let x represent the decision variables associated with the design of the locker network and V(x) denote the set of locations where lockers are installed. 
% Assume that Z ̃ N ̃ (x) corresponds to the locations of home deliveries that remain after accounting for locker installations. 
% Our research focuses on developing models for two distinct operational systems: 
% • Two-trip mode: This mode is modeled as the sum of two separate traveling salesman problems (TSPs):  TSP(V(x)) + TSP(Z ̃ N ̃ (x)): 
% Here, the first term represents the TSP for the locker  locations, and the second term represents the TSP for the remaining home delivery locations. 
% • Mixed-trip mode: This mode is modeled as a single TSP that combines both locker locations and remaining home delivery locations:  TSP(V(x) ∪ Z ̃ N ̃ (x)):  This unified approach reflects a scenario in which deliveries to lockers and homes are intermingled in a single route.

% Contributions
\color{black}
We contribute as follows:
\begin{itemize}
    \item \textit{An integrated demand-responsive network design formulation}:
    We connect a user-choice demand interface to the joint selection of physical drone bases, \textcolor{blue}{unified pickup/drop-off smart lockers}, battery-swapping stations, \textcolor{blue}{smart-locker capacities}, flight connections, and fleet size.
    The interface allows infrastructure decisions to change the displayed locker alternative and the fulfillment-cost composition of successful service without introducing backend operational variables into customer utility; in the locker-only restriction, it also determines the demand offered to the drone network.
    \item \textit{A tractable representation of availability, chaining, and congestion}:
    We use discrete availability levels and continuous offered-flow allocations to represent endogenous task assignment across physical bases and virtual idle-drone states.
    One closed-network flow-balance system links realized demand, logical state transitions, empty repositioning, \textcolor{blue}{smart-locker congestion}, and fleet requirements in a MISOCP.
    \item \textit{A planning framework for operational and managerial evaluation}:
    The formulation supports experiments on infrastructure mix, user adoption, service reliability, battery range, task chaining, network structure, and drone utilization.
    These experiments are designed to quantify the trade-offs among courier-cost savings, infrastructure investment, service coverage, and fleet efficiency.
    \textcolor{blue}{The resulting evidence characterizes cost and service trade-offs relevant to commercial viability but does not evaluate profitability itself.}

\end{itemize}

% Outline
The remainder of this paper is organized as follows.
Section~\ref{sec:literature_review} reviews the related literature.
Section~\ref{sec:formulation} describes the planning problem and presents the integrated formulation.
Section~\ref{sec:solution_methods} records candidate network valid inequalities for later computational development.
Section~\ref{sec:experiments} presents the computational study and managerial analyses.
Section~\ref{sec:conclusion} concludes the paper.


\section{Literature Review}
\label{sec:literature_review}

Our research intersects three main streams of literature within operations research and supply chain management: network design for low-altitude logistics, hierarchical facility location with a focus on charging infrastructure, and facility location models incorporating congestion and queueing effects. 

\subsection{Network Design for Low-Altitude Logistics}

The integration of unmanned aerial vehicles (UAVs), or drones, into commercial and emergency logistics has spurred a new domain of network design focused on low-altitude operations.
Due to the limited flight range of battery-powered drones, multiple recharging stations are required in large urban areas to complete deliveries without running out of power \citep{hongDeviationFlow2017AdvancesinGeocomputation}.
Consequently, the optimal placement of these bases and the routing of drones heavily depend on continuous space shortest paths and range-restricted flow-refueling logic \citep{hongRangerestrictedRechargingStation2018}.

Recently, drone networks have been increasingly deployed to deliver time-critical materials, such as medication, blood, and exam kits, in emergency situations \citep{wangFacilityLocation2022}.
In this context, \citet{boutilierDroneNetwork2022M&SOM} developed an integrated location-queuing model based on a $p$-median architecture to design a defibrillator-enabled drone network that augments existing emergency medical services.
Similarly, \citet{lejeuneDroneDeliveryNetwork2025OperationsResearch} proposed a stochastic emergency network design model for opioid overdose response, formulating the network as a collection of M/G/K queueing systems where the capacity of each system is an endogenous decision variable.
Because the resulting nonlinear integer problem is intractable in its original form, they developed a mixed-integer linear programming (MILP) reformulation and a branch-and-cut algorithmic framework.

To account for system congestion in drone-based emergency delivery, \citet{wangFacilityLocation2022} jointly optimized the location of facilities, the capacity of drones at opened facilities, and the allocation of demands.
They employed queues to model the system congestion of drone requests under non-priority, static priority, and dynamic priority disciplines, subsequently approximating the model as a computationally tractable mixed-integer second-order conic program (MISOCP).

Beyond emergency logistics, \citet{chenCoordinatedDroneCourier2025} recently introduced a coordinated drone-courier logistics system for intra-city express services, recognizing that drones are often unable to execute ``last-meter'' deliveries directly to customer doorsteps due to airspace and safety restrictions.
In their system, drones operate within a closed network among vertiports, while couriers connect customers to the drone delivery system.
They formulated an integrated optimization model where the interplay between network planning and tactical operations is mirrored by a queueing network model.
To handle the resulting nonlinear and partially convex/concave feasible region, they developed an iterative exact algorithm that adaptively refines linear approximations of the nonlinear constraints.
Our work extends this stream by explicitly capturing chaining tasks and unified energy balances across a multi-type drone network.

\color{blue}
Alongside operational deployment, a growing literature examines the economic conditions under which drone delivery can compete with ground service.
% Citation link for \citep{DroneEconomicViabilityReview2025}:
% https://doi.org/10.1108/IJLM-04-2023-0149
Existing evidence identifies interdependent technological, operational, economic, regulatory, and infrastructure conditions that shape drone-delivery performance and adoption \citep{DroneEconomicViabilityReview2025}.
Accordingly, we distinguish demonstrated operational viability from unresolved commercial viability and interpret modeled fulfillment-cost outcomes only as evidence on the cost dimension of commercial viability, not as proof of profitability.
\color{black}

\subsection{Hierarchical Facility Location and Charging Infrastructure}

The second stream of relevant literature encompasses hierarchical facility location problems (HFLP) and the optimal deployment of alternative fuel and charging infrastructure.
The primary objective in a typical hierarchical facility location problem is to determine the location of facilities in a multi-level network in a way to serve the customers at the lowest level of the hierarchy both efficiently and effectively \citep{farahaniHierarchicalFacilityLocation2014}.

With the electrification of transportation, a significant portion of this literature has pivoted toward the charging station location problem (CSLP), which seeks to optimize the locations of charging stations to satisfy customer recharging demand under real-life restrictions \citep{kchaou-boujelbenChargingStation2021TRC}.
A widely adopted paradigm is the flow-refueling location model (FRLM), which positions refueling facilities to maximize the flow that can be accommodated in a road network while respecting the range limitations of vehicles \citep{mirhassaniFlexibleReformulationRefueling2013}.
Extensions of the FRLM have explicitly added the routing aspect of individual drivers and their shortest path deviation tolerances \citep{yildizBranchPriceApproach2016}.

However, charging an electric vehicle (EV) is inherently a time-consuming burden, which can lead to severe congestion if the public charging infrastructure is not well-designed \citep{liRobustPlanning2025TransportationResearchPartB:Methodological}.
Addressing this, \citet{liRobustPlanning2025TransportationResearchPartB:Methodological} studied the strategic planning of public EV charging stations under uncertain charging demand and service rates, introducing chance constraints to mitigate long waiting times.

A major paradigm shift in mitigating charging times is the joint deployment of standard chargers and battery swapping stations.
\citet{laiMultimodalCharging2024TransportationResearchPartB:Methodological} proposed a multimodal charging network in which charging stations and battery swapping stations are jointly built to synergistically support an electric ride-hailing fleet.
Their rationale is that while charging stations are cost-effective for scaling up fleets, battery swapping stations offer quick turnaround and can improve fleet utilization \citep{laiMultimodalCharging2024TransportationResearchPartB:Methodological}.
Incorporating elastic passenger demand and spatial charging equilibrium, they formulated the problem as a nonconcave program and established tight upper bounds through relaxation and decomposition \citep{laiMultimodalCharging2024TransportationResearchPartB:Methodological}.

Furthermore, the operational integration of charging infrastructure within vehicle sharing systems is heavily explored.
\citet{heChargingElectric2020M&SOM} investigated how to charge an EV-sharing fleet to ensure viability and profitability by integrating charging infrastructure planning and vehicle repositioning operations.
They developed a queuing network model that characterizes how customers endogenously pick EVs based on energy levels and how the operator implements charging policies.
Similarly, \citet{heServiceRegion2017M&SOM} formulated a mathematical programming model for service region design that incorporates customer adoption behavior and fleet management, addressing inherent planning uncertainty through a distributionally robust optimization framework approximated by an MISOCP.
Our model draws inspiration from these hierarchical and operational frameworks but applies them to the unique requirements of drone base and battery swapping station co-location.

\subsection{Facility Location with Congestion and Queueing Effects}

The traditional facility location literature often assumes infinite processing capacity at facilities, ignoring the waiting times incurred by users.
However, in service networks, the facility location problem with immobile servers, stochastic demand, and congestion is highly critical, often modeled as a network of M/M/1 queues leading to complex nonlinear mixed-integer programs \citep{elhedhliServiceSystem2006M&SOM}.

To model the flow of entities in constrained logistics networks, researchers have utilized closed queueing networks.
For instance, \citet{adelmanPriceDirectedControl2007OperationsResearch} considered an internal pricing mechanism for managing a fleet of service units flowing in a closed logistics queueing network, utilizing Palm calculus to approximate queueing and stockout effects.
\citet{shuModelsEffective2013OperationsResearch} developed a network flow model with proportionality constraints to estimate the flow of bicycles within a public bicycle-sharing system, given an initial allocation of bicycles.
\citet{heServiceRegion2017M&SOM} developed closed queue network for electric vehicle sharing systems, and the ideas are extended to air mobility network \citep{wangVertiportPlanning2022M&SOM,wangUrbanAir2025TRE} and intra-city express services \citep{chenCoordinatedDroneCourier2025}.

The primary mathematical challenge in these congestion models lies in handling the nonlinear, fractional queueing formulas within integer programming frameworks.
While earlier works relied on piecewise linear approximations \citep{elhedhliServiceSystem2006M&SOM}, recent methodological advancements have proven that mixed-integer second-order cone programs (MISOCPs) form a highly efficient class of mixed-integer convex programs for these problems \citep{ahmadi-javidConvexificationQueueing2022INFORMSJournalonComputing}.
\citet{ahmadi-javidConvexificationQueueing2022INFORMSJournalonComputing} comprehensively demonstrated how various performance metrics of M/G/1 queues can be modeled by different MISOCPs.
They showed that conic reformulation methods efficiently solve very large-size test problems that cannot be solved by existing methods based on linear programming or state-space decomposition within reasonable times.

This SOCP convexification approach has become the state-of-the-art technique in recent literature.
For example, \citet{heChargingElectric2020M&SOM} utilized lower and upper bound formulations as mixed-integer second-order cone programs to achieve computationally tractable solutions for their integrated queuing-location EV model.
\citet{wangFacilityLocation2022} similarly leveraged MISOCP reformulations to solve their drone facility location model with congestion and priority queuing.
Aligning with this advanced methodological stream, our work reformulates the highly non-linear fractional terms associated with drone waiting times and base congestion into an SOCP.
This allows us to rigorously incorporate congestion effects and chaining logic while maintaining computational tractability using off-the-shelf solvers.


\section{Problem Description and Assumptions}
\label{sec:assumptions}


\textcolor{blue}{Motivated by the demonstrated operational viability of drone delivery in selected applications, this section shifts attention to the cost-related planning conditions that may support, but do not by themselves establish, commercial viability.}
Planning a supplementary urban drone delivery network requires infrastructure design and its long-run operating consequences to be evaluated jointly. Opening a customer-facing smart locker changes the delivery alternatives available to customers and the fulfillment-cost composition of successful orders; in the locker-only restriction, it also determines the order volume eligible for the drone network. The volume selected for drone service determines drone movements, which in turn determine vehicle availability, locker congestion, energy-feasible routing, and fleet requirements. These interactions are further complicated by stochastic order arrivals and the circulation of a finite drone fleet. This section therefore describes the common service architecture, distinguishes the general mixed-delivery problem from its two service-scope restrictions, and states the assumptions that connect customer choice, courier fulfillment, and drone operations.

Figure 2 summarizes the physical drone movements shared by all three service scopes. For any order admitted to drone service, an idle drone is dispatched from a physical base to an origin smart locker or released from a virtual idle state associated with that locker. A courier brings the parcel from its origin to the origin locker, after which the drone transports it to a destination locker. At the destination locker, a locker-pickup order awaits customer collection, whereas a home-delivery order is transferred to a courier for the final leg to the customer. After unloading, the drone either returns empty to a physical base or enters a virtual idle state near a locker for a follow-on assignment. The solid links in Figure 2 represent direct physical movements, whereas the dashed links represent movements that use at most one battery-swapping station. Virtual idle states are operational constructs and are therefore not depicted as physical facilities.

\color{blue}
The three service scopes differ only in the order classes that may enter this common drone network.
The general mixed-delivery problem allows locker-pickup orders and platform-selected home-delivery orders to share the drone network without a predetermined priority.
The home-delivery-assistance restriction retains only drone-assisted home delivery, whereas the locker-only restriction admits only choice-generated potential locker-pickup demand to drone service.
The term locker-only therefore describes the scope of drone service rather than the platform's full service portfolio.

The customer-choice layer estimates the potential preference for locker pickup but does not commit every such order to locker fulfillment or determine the platform's internal transport plan.
Home delivery remains available in all three service scopes, whereas the potential locker alternative may be configured in the general problem and the locker-only restriction.
When both alternatives are considered, the choice model estimates potential locker uptake from delivery time, price, and walking convenience, and the platform may display one of the nearest eligible open lockers.
The potential locker share becomes completed locker pickup only for the corresponding flow realized by the drone network; all remaining demand is fulfilled by direct courier home delivery.
The home-delivery alternative is evaluated using an exogenous scenario delivery-time attribute rather than an estimate from the order data, as detailed in Section~\ref{subsec:data_instances}.
A home-delivery customer does not observe whether the platform subsequently assigns the middle-mile movement to drones or which internal locker pair is used.
These backend decisions therefore do not enter customer utility.
Delivery prices and convenience attributes are treated as exogenous inputs to customer choice rather than as platform decisions, and the model contains no service-revenue term.
\color{black}

\color{blue}
The courier fallback rule guarantees fulfillment when drone capacity or availability is insufficient.
A home-delivery order that is not assigned to or ultimately served by the drone network is transported directly from its origin to its destination by couriers.
Likewise, potential locker demand that is not realized by the drone network remains in the home-delivery alternative and is transported directly from its origin to its customer destination by a courier; it never becomes a committed locker-pickup order.
Courier capacity and courier congestion are not modeled explicitly.
Instead, the existing courier system is assumed to have sufficient capacity to absorb all residual demand at exogenous fulfillment costs, consistent with the role of drones as a supplementary rather than stand-alone delivery resource.
This fallback assumption supports strategic cost comparisons but excludes endogenous courier wages, peak-capacity shortages, and courier congestion from the commercial interpretation.
\color{black}

The infrastructure layer determines where and how the drone network can operate. The platform may open three types of physical facilities: drone bases that store, charge, inspect, and dispatch drones; customer-facing smart lockers that serve as pickup, drop-off, and transfer points; and battery-swapping stations that extend the feasible range of physical movements. The platform also selects smart-locker handling capacities, physical flight connections, and the drone fleet size. Virtual idle states support task chaining after a delivery, but they are bookkeeping states rather than physical facilities and therefore require neither siting nor physical connection decisions.

The operational layer captures the capacity bottlenecks that are material to strategic planning while abstracting from real-time dispatch details. Orders in each origin-destination market are assumed to arrive according to a Poisson process, and drones, rather than parcels, circulate in a closed queueing network. At any time, each drone is idle at a physical base or in a virtual state, undergoing handling at a smart locker, or traveling on a physical movement. Smart-locker handling is represented as a single-server queue, whereas flight movements are represented as infinite-server transit states. Because every drone-served order is processed at an origin and a destination locker, limited locker handling capacity can create congestion. The model therefore constrains average locker-handling and virtual-state waiting times. For tractability, the model uses long-run flow rates and discrete availability levels, screens each direct or one-swap route for route-level battery feasibility, and imposes aggregate energy sufficiency instead of tracking the state of charge of each drone.

% Previous wording retained for revision history: The planning objective is to identify the least-cost feasible system design and fulfillment plan for each operating mode. The model jointly determines physical facility locations, smart-locker capacities, network connections, fleet size, customer-facing locker service configurations, and the allocation of eligible home-delivery orders to the drone network. The objective accounts for infrastructure, connection, flight, fleet, locker-handling, and courier-fulfillment costs. Because courier fallback ensures that all demand is fulfilled, the model does not require every order to be served by drones. Instead, it determines when drone service is cost-effective while respecting flow-balance, service-quality, and battery-feasibility requirements.
\color{blue}
For the general problem and each restriction, the planning objective identifies the least-cost design and fulfillment plan that satisfies the model's capacity, congestion, flow-balance, service-quality, and battery-feasibility requirements.
Across these problems, the platform determines the common infrastructure, capacity, connection, and fleet decisions and, where applicable, customer-facing locker configurations and home-delivery allocations.
Its objective includes the explicitly specified infrastructure, connection, flight, fleet, locker-service, and courier-fulfillment costs.
Because courier fallback fulfills all residual demand, the model does not require every order to be served by drones.
At the system level, the optimized solution can use drone service only when doing so does not increase modeled platform fulfillment cost relative to an available courier-fallback solution under the stated assumptions.
Whether drone service produces a strict cost saving is determined through explicit counterfactual comparisons in the computational study.
A strict modeled cost saving provides evidence on one cost dimension relevant to commercial viability and should not be interpreted as evidence of profitability.


\begin{figure*}[!t]
\centering
\begin{tikzpicture}
\node[anchor=south west,inner sep=0] (networkfigure) at (0,0)
{\includegraphics[width=5.5in]{Figure/Intro.png}};
\begin{scope}[x={(networkfigure.south east)},y={(networkfigure.north west)}]
\fill[white] (0.000,0.238) rectangle (0.168,0.350);
\node[text=blue,font=\scriptsize\sffamily] at (0.078,0.294)
{Drone Base ($\mathcal L$)};
\fill[white] (0.272,0.158) rectangle (0.410,0.270);
\node[text=blue,font=\scriptsize\sffamily] at (0.336,0.214)
{Smart Lockers ($\mathcal H$)};
\fill[white] (0.574,0.238) rectangle (0.724,0.350);
\node[text=blue,font=\scriptsize\sffamily] at (0.642,0.294)
{Smart Lockers ($\mathcal H$)};
\fill[white] (0.826,0.238) rectangle (1.000,0.350);
\node[text=blue,font=\scriptsize\sffamily] at (0.919,0.294)
{Drone Base ($\mathcal L$)};
\end{scope}
\end{tikzpicture}
\caption{Overview of the supplementary drone delivery network.}
\label{fig:intro_network}
\vspace{-0.10in}
\end{figure*}

% 首段需要回答：为什么要用这么复杂的模型？因为场景很多，随机性很多，然后说明这一章要干嘛

% 在合适的地方加入配图的具体运营流程，并且需要把运营流程，我们的假设和三种模式写的更清楚

% 现在的问题是，下面写的这些都是mixed，好像我们默认了讲的是这种模式。这里需要弄明白是以一个举例还是以一种通用的方式，不然写的太乱了

% Previous draft retained for revision history: From the customer's perspective, the platform offers two mutually exclusive delivery options: home delivery and locker pickup.
% Previous draft retained for revision history: Customers' preferences are related to the delivery time, price, and convenience of each option.
% Previous draft retained for revision history: For those home-delivery orders, the platform may selectively offload the middle-mile movement to the drone network, while the remaining portion is fulfilled by human couriers.
% Previous draft retained for revision history: However, we assume this process is invisible to the customer, who only observes the displayed estimated home-delivery time of human couriers by empirical data.
% Previous draft retained for revision history: If the infrastructure is not configured to offer locker pickup, all potential locker-pickup orders remain in the home-delivery market.
% Previous draft retained for revision history: Once locker pickup is activated, any order not realized by the drone network because of stochastic drone unavailability is fulfilled by couriers directly to the customer destination as home delivery.
% Previous draft retained for revision history: On the operation side, the platform may selectively offload home-delivery demand to the drone network, while activated locker demand is offered to drones.
% Previous draft retained for revision history: The portions of either order class that are not realized by drones are fulfilled by human couriers.
% Previous draft retained for revision history: However, the courier system is not explicitly modeled, the human couriers are assumed to be available always, which is consistent to the intuition that, drone delivery network is a potential complementary resources service to the existing courier system.
% Previous draft retained for revision history: In the aspect of infrastructure side, three types of candidate facilities are considered: physical drone bases, customer-facing smart lockers, and battery-swapping stations.
% Previous draft retained for revision history: Due to the limited availabilities of drones, drone energy capacity, locker cpacities. The congestions effect of the customer-facing lockers are significant. Which could be one of the main bottleneck of the drone delivery network.
% Previous draft retained for revision history: And the platform's objective is to minimize the total cost of drone operations and infrastructure investment while satisfying customer demand.
The platform therefore treats drones as a supplementary resource and relies on the existing courier system to fulfill residual demand.
The joint planning problem coordinates physical drone bases, customer-facing smart lockers, battery-swapping stations, drone availability, and locker congestion under the general formulation and its two restrictions.
The platform minimizes modeled fulfillment cost, including the represented drone-network, infrastructure, locker-service, and courier-fallback costs, while fulfilling all demand.

\color{black}


\section{Model Formulation}
\label{sec:path_model}

We consider a meal-delivery platform that combines couriers and drones to serve origin--destination (OD) demand.
The platform jointly selects facilities, locker capacities, and flight links, assigns each OD pair to a departure--return base pair, and determines drone availability and path flows.
The model integrates these decisions through a closed queueing network and complete drone service routes.
All demand rates, service times, and costs are expressed in consistent time units.

\subsection{Strategic Network Decisions}
\label{subsec:path_design}

Let $\mathcal D$ be the set of OD pairs, where $o$ denotes a merchant and $d$ denotes a customer location.
The potential order rate for $(o,d)$ is $h_{o,d}>0$.
Let $\mathcal L$, $\mathcal H$, and $\mathcal E$ denote the candidate drone bases, locker locations, and independent battery-swapping stations, respectively.
The complete facility set is $\mathcal V=\mathcal L\cup\mathcal H\cup\mathcal E$, and batteries can be replaced at nodes in $\mathcal K=\mathcal L\cup\mathcal E$.
Bases are indexed by $\ell,m$, locker locations by $i,j$, and replenishment nodes by $k$.

The binary variable $x_v$ indicates whether facility $v\in\mathcal V$ is opened.
Locker capacity is selected from a finite set of installation levels $\mathcal S$.
The binary variable $z_{i,s}$ equals one if level $s$ is selected at location $i$; this level provides handling capacity $\mu_{i,s}$ and incurs fixed cost $f_{i,s}^{\mathrm{locker}}$.
Exactly one level is selected at an open locker location.
For identical cabinets and $\mathcal S=\{1,\ldots,5\}$, level $s$ installs $s$ cabinets and provides $\mu_{i,s}=s\mu_i$, where $\mu_i$ is the handling capacity of one cabinet.
The fixed costs of a base and an independent swapping station are $f_\ell^{\mathrm{base}}$ and $f_k^{\mathrm{swap}}$, respectively.

The binary variable $\alpha_{o,\ell,m,d}$ assigns OD pair $(o,d)$ to departure base $\ell$ and return base $m$.
Each OD pair is assigned to at most one pair of open bases, and $\ell=m$ is allowed.
Orders assigned to a base pair can be split continuously among its feasible service routes.
Orders without a base assignment are delivered directly by couriers.

Let $\mathcal A$ be the set of candidate directed flight links between facility nodes.
The binary variable $y_{u,v}$ indicates whether link $(u,v)\in\mathcal A$ is enabled, at fixed cost $f_{u,v}^{\mathrm{link}}$.
Both endpoint facilities must be open, and the two directions are selected separately.
When distinct facility functions share a location, their zero-time connections form the subset $\mathcal A^0\subseteq\mathcal A$.
Such a connection has zero cost and is enabled whenever both endpoint functions are open.

\subsection{Drone Operations}
\label{subsec:path_operations}

\subsubsection{Queueing Network}
\label{subsubsec:path_queue}

An order requests a drone when its courier collects the meal from the merchant.
We assume that these requests form independent stationary Poisson streams with rates $h_{o,d}$.
If a drone is available at the assigned departure base, the platform selects a service route and dispatches the drone toward its pickup locker.
Otherwise, the courier delivers the order directly to the customer.
Meal preparation and the courier's journey to the merchant precede the drone request.

A drone completes its service route at the assigned return base after replenishing its battery.
At that point, the drone either joins the local idle inventory or is immediately repositioned to another base.
The same decision is made after an empty repositioning route is completed.
These routing decisions are independent of instantaneous inventory and use stationary probabilities determined by the planned flows.
Consequently, repositioning acts on returning drones before they enter an idle queue; a base can forward arriving drones while holding no idle inventory.

The queueing network contains single-server idle nodes at bases and infinite-server activities for complete service and repositioning routes.
Idle service completions correspond to order requests acquiring drones, while route activities hold drones until their return and replenishment are complete.
Returning activities can feed either an idle node or a subsequent repositioning activity.
Following the fixed-population-mean (FPM) planning approach of \citet{heServiceRegion2017M&SOM} and Wang et al.~\citeyearpar{wangVertiportPlanning2022M&SOM}, we use $\rho_\ell$ for the availability of a drone at base $\ell$.

Let $\mathcal P_{o,\ell,m,d}$ be the set of feasible service paths for OD pair $(o,d)$ and base pair $(\ell,m)$.
The actual order rate on path $p\in\mathcal P_{o,\ell,m,d}$ is $\psi_{\ell,m}^{p}\ge0$.
Each path retains its OD identity through its path set.
The successful service rate therefore satisfies
\begin{equation}
 \sum_{p\in\mathcal P_{o,\ell,m,d}}\psi_{\ell,m}^{p}
 =h_{o,d}\alpha_{o,\ell,m,d}\rho_\ell.
 \label{eq:path_product}
\end{equation}
Because $\alpha$ is binary and each OD pair selects at most one base pair, \eqref{eq:path_product} has the following exact linear representation:
\begin{subequations}\label{eq:path_linearization}
\begin{align}
 \sum_{p\in\mathcal P_{o,\ell,m,d}}\psi_{\ell,m}^{p}
 &\le h_{o,d}\alpha_{o,\ell,m,d},
 &&(o,d)\in\mathcal D,\ \ell,m\in\mathcal L,\label{eq:path_pair}\\
 \sum_{m\in\mathcal L}\sum_{p\in\mathcal P_{o,\ell,m,d}}\psi_{\ell,m}^{p}
 &\le h_{o,d}\rho_\ell,
 &&(o,d)\in\mathcal D,\ \ell\in\mathcal L,\label{eq:path_source_upper}\\
 h_{o,d}\rho_\ell-
 \sum_{m\in\mathcal L}\sum_{p\in\mathcal P_{o,\ell,m,d}}\psi_{\ell,m}^{p}
 &\le h_{o,d}\left(1-\sum_{m\in\mathcal L}\alpha_{o,\ell,m,d}\right),
 &&(o,d)\in\mathcal D,\ \ell\in\mathcal L.\label{eq:path_source_lower}
\end{align}
\end{subequations}
At the selected departure base, the two source inequalities fix total service at $h_{o,d}\rho_\ell$, and the pair inequality assigns that flow to the selected return base.
At every unselected departure base, the pair inequalities set all corresponding service flows to zero.

Let $\mathcal R_{\ell,m}$ be the feasible empty repositioning paths from $\ell$ to $m$, for $\ell\ne m$, with actual rates $\varphi_{\ell,m}^{r}\ge0$.
The total service and repositioning flow entering each base equals its total outgoing flow.
These balances also determine the stationary routing probabilities at replenished returns.

Let $I_\ell$ be the mean number of idle drones at base $\ell$.
The FPM idle-inventory relationship is $I_\ell\ge\rho_\ell/(1-\rho_\ell)$.
Its equivalent second-order cone representation is
\begin{equation}
 \left\|\begin{pmatrix}2\\I_\ell+\rho_\ell\end{pmatrix}\right\|_2
 \le I_\ell+2-\rho_\ell,
 \qquad I_\ell\ge0,\qquad 0\le\rho_\ell\le x_\ell.
 \label{eq:path_idle_cone}
\end{equation}
For finite $I_\ell$, this cone enforces $\rho_\ell<1$.
With $\tau_p$ and $\tau_r$ denoting mean service and repositioning durations, Little's law gives the fleet requirement
\begin{equation}
 \sum_{\ell\in\mathcal L}I_\ell
 +\sum_{(o,d)\in\mathcal D}\sum_{\ell,m\in\mathcal L}
   \sum_{p\in\mathcal P_{o,\ell,m,d}}\tau_p\psi_{\ell,m}^{p}
 +\sum_{\ell\ne m}\sum_{r\in\mathcal R_{\ell,m}}\tau_r\varphi_{\ell,m}^{r}
 \le\Gamma,
 \label{eq:path_fleet}
\end{equation}
where $\Gamma\ge0$ is a continuous fleet-sizing decision with unit cost $f^{\mathrm{fleet}}>0$.
This relationship combines idle inventory with drones engaged in complete routes, as in the fleet-planning formulations of \citet{heServiceRegion2017M&SOM,chenCoordinatedDroneCourier2025}.

\subsubsection{Drone Routes and Service Times}
\label{subsubsec:path_times}

A service path specifies a pickup locker $i(p)$, a delivery locker $j(p)\ne i(p)$, and the three successive stages
\begin{equation}
 \ell\longrightarrow i(p),\qquad i(p)\longrightarrow j(p),\qquad j(p)\longrightarrow m.
 \label{eq:path_stages}
\end{equation}
Intermediate replenishment is allowed at nodes in $\mathcal K$.
The only locker visits on a service path are its designated pickup and delivery actions, and all intermediate nodes belong to $\mathcal K$.
An empty repositioning path visits replenishment nodes only.
Let $t_{u,v}^{\mathrm{fly}}$ be the flight time on a candidate link, $t^{\mathrm{bat}}$ the usable flight-time budget of a full battery, and $t_k^{\mathrm{sw}}$ the replacement time at node $k$.
Energy consumption is proportional to flight time and is the same for loaded and empty flight.
The accumulated flight time between successive replenishments cannot exceed $t^{\mathrm{bat}}$.
The drone starts fully replenished, preserves its residual energy at lockers, and receives a full replacement at every visited replenishment node, including its final return base.
An intermediate base visit provides replenishment while the drone remains committed to its route.
Replacement resources are sufficient, and a drone waiting for a courier rests on the ground.

Let $t_p^1$ be the time from departure to arrival at the pickup locker, $t_p^2$ the time from meal collection there to deposition at the delivery locker, and $t_p^3$ the subsequent return-and-replenishment time.
Each stage includes its flight, handling, and intermediate replacement times.
The courier's travel time from merchant $o$ to pickup locker $i$ is the nonnegative random variable $T_{o,i}$.
Route selection precedes its realization, and courier times are independent across orders and independent of inventory at the request epoch.
We use an exponential specification with mean $t_{o,i}^{\mathrm r}$, estimated from delivery data.
The complete service duration and its mean are
\begin{align}
 T_p&=\max\{t_p^1,T_{o,i(p)}\}+t_p^2+t_p^3,\label{eq:path_random_time}\\
 \tau_p&=t_p^1+t_p^2+t_p^3
 +t_{o,i(p)}^{\mathrm r}\exp\left(-\frac{t_p^1}{t_{o,i(p)}^{\mathrm r}}\right).
 \label{eq:path_mean_time}
\end{align}
If $t_{o,i(p)}^{\mathrm r}=0$, the last term is zero.
These durations define the general service-time distributions of the infinite-server activities; the planning model uses their means.
An empty path has duration $\tau_r$, including all flight and replacement times.

For each path, $a_{u,v}^{p}$ or $a_{u,v}^{r}$ counts its traversals of directed link $(u,v)$, and $t_p^{\mathrm{fly}}$ or $t_r^{\mathrm{fly}}$ is its total flight time.
The locker coefficient $b_i^p=\mathbf1\{i(p)=i\}+\mathbf1\{j(p)=i\}$ records the handling event at location $i$.
Path sets are constructed on the full candidate graph and incorporate the energy and action-order requirements.
Within each stage, a replenishment node need only be visited once; an empty route likewise need not repeat a replenishment node.
Removing a repeated-node cycle preserves the full battery state and weakly reduces time and cost.

\subsection{Delivery-Mode Choice}
\label{subsec:path_choice}

At destination locker $j$, a customer at $d$ chooses between self-pickup and home delivery by a courier.
Their deterministic utility components are $V_{d,j}^{P}$ and $V_{d,j}^{H}$, based on the specified walking, price, and terminal-service attributes.
With independent type-I extreme-value utility disturbances, the self-pickup probability is
\begin{equation}
 \sigma_{d,j}=\frac{\exp(V_{d,j}^{P})}
 {\exp(V_{d,j}^{P})+\exp(V_{d,j}^{H})}.
 \label{eq:path_choice}
\end{equation}
Self-pickup is excluded when the locker is outside the allowed collection range, in which case $\sigma_{d,j}=0$.
The utility attributes are fixed for each customer--locker pair, so the probabilities are evaluated before optimization.
The path allocation selects the destination locker and therefore the applicable choice probability.

Let $c_{a,b}^{\mathrm{rider}}$ be the per-order courier cost from location $a$ to $b$, and let $c^{\mathrm{drone}}$ be the drone cost per unit of flight time.
Self-pickup incurs zero terminal delivery cost, while home delivery incurs $c_{j,d}^{\mathrm{rider}}$.
The operating costs of service and repositioning columns are consequently
\begin{align}
 c_p^{\mathrm{service}}
 &=c_{o,i(p)}^{\mathrm{rider}}
 +(1-\sigma_{d,j(p)})c_{j(p),d}^{\mathrm{rider}}
 +c^{\mathrm{drone}}t_p^{\mathrm{fly}},
 &&p\in\mathcal P_{o,\ell,m,d},\label{eq:path_service_cost}\\
 c_r^{\mathrm{reposition}}
 &=c^{\mathrm{drone}}t_r^{\mathrm{fly}},
 &&r\in\mathcal R_{\ell,m}.\label{eq:path_empty_cost}
\end{align}
Fleet costs are charged separately through $\Gamma$.

\subsection{Full Formulation}
\label{subsec:path_full}

The objective minimizes facility and flight-link costs, fleet costs, path operating costs, and the direct-courier cost of all orders not served by drones.
Each realized drone-assisted order removes one order from the direct-courier flow.
For link activation, define $M^0=(|\mathcal L|+2)\sum_{(o,d)\in\mathcal D}h_{o,d}$ and
\begin{equation}
 M_{u,v}=\min\left(\{M^0\}\cup
 \left\{\max_{s\in\mathcal S}\mu_{i,s}:i\in\{u,v\}\cap\mathcal H\right\}\right).
 \label{eq:path_link_bound}
\end{equation}
An optimal operating plan can be chosen to satisfy these bounds; Section~\ref{subsec:path_inequalities} establishes them using empty-cycle removal and locker capacities.
The complete mixed-integer second-order cone formulation is
\begin{subequations}\label{eq:path_full}
\begin{align}
 \min\quad
 &\sum_{\ell\in\mathcal L}f_\ell^{\mathrm{base}}x_\ell
 +\sum_{i\in\mathcal H}\sum_{s\in\mathcal S}f_{i,s}^{\mathrm{locker}}z_{i,s}
 +\sum_{k\in\mathcal E}f_k^{\mathrm{swap}}x_k
 +\sum_{(u,v)\in\mathcal A}f_{u,v}^{\mathrm{link}}y_{u,v}
 +f^{\mathrm{fleet}}\Gamma\notag\\
 &+\sum_{(o,d)\in\mathcal D}\sum_{\ell,m\in\mathcal L}
   \sum_{p\in\mathcal P_{o,\ell,m,d}}c_p^{\mathrm{service}}\psi_{\ell,m}^{p}
 +\sum_{\ell\ne m}\sum_{r\in\mathcal R_{\ell,m}}
   c_r^{\mathrm{reposition}}\varphi_{\ell,m}^{r}\notag\\
 &+\sum_{(o,d)\in\mathcal D}c_{o,d}^{\mathrm{rider}}
  \left(h_{o,d}-\sum_{\ell,m\in\mathcal L}
  \sum_{p\in\mathcal P_{o,\ell,m,d}}\psi_{\ell,m}^{p}\right),
 \label{eq:path_objective}\\
 \text{s.t.}\quad
 &\sum_{s\in\mathcal S}z_{i,s}=x_i,
 &&i\in\mathcal H,\label{eq:path_install}\\
 &\sum_{\ell,m\in\mathcal L}\alpha_{o,\ell,m,d}\le1,
 &&(o,d)\in\mathcal D,\label{eq:path_assign}\\
 &\sum_{m\in\mathcal L}\alpha_{o,\ell,m,d}\le x_\ell,
 &&(o,d)\in\mathcal D,\ \ell\in\mathcal L,\label{eq:path_open_source}\\
 &\sum_{\ell\in\mathcal L}\alpha_{o,\ell,m,d}\le x_m,
 &&(o,d)\in\mathcal D,\ m\in\mathcal L,\label{eq:path_open_return}\\
 &y_{u,v}\le x_u,\quad y_{u,v}\le x_v,
 &&(u,v)\in\mathcal A,\label{eq:path_open_link}\\
 &y_{u,v}\ge x_u+x_v-1,
 &&(u,v)\in\mathcal A^0,\label{eq:path_colocation}\\
 &0\le\rho_\ell\le x_\ell,
 &&\ell\in\mathcal L,\label{eq:path_rho}\\
 &\|(2,I_\ell+\rho_\ell)\|_2\le I_\ell+2-\rho_\ell,
 &&\ell\in\mathcal L,\label{eq:path_full_cone}\\
 &\sum_{p\in\mathcal P_{o,\ell,m,d}}\psi_{\ell,m}^{p}
 \le h_{o,d}\alpha_{o,\ell,m,d},
 &&\substack{(o,d)\in\mathcal D,\\\ell,m\in\mathcal L},\label{eq:path_full_pair}\\
 &\sum_{m\in\mathcal L}\sum_{p\in\mathcal P_{o,\ell,m,d}}\psi_{\ell,m}^{p}
 \le h_{o,d}\rho_\ell,
 &&(o,d)\in\mathcal D,\ \ell\in\mathcal L,\label{eq:path_full_source_upper}\\
 &h_{o,d}\rho_\ell-
 \sum_{m\in\mathcal L}\sum_{p\in\mathcal P_{o,\ell,m,d}}\psi_{\ell,m}^{p}
 \le h_{o,d}\left(1-\sum_{m\in\mathcal L}\alpha_{o,\ell,m,d}\right),
 &&(o,d)\in\mathcal D,\ \ell\in\mathcal L,\label{eq:path_full_source_lower}\\
 &\sum_{(o,d)\in\mathcal D}\sum_{m\in\mathcal L}
  \sum_{p\in\mathcal P_{o,\ell,m,d}}\psi_{\ell,m}^{p}
 +\sum_{m\ne\ell}\sum_{r\in\mathcal R_{\ell,m}}\varphi_{\ell,m}^{r}\notag\\
 &\qquad=\sum_{(o,d)\in\mathcal D}\sum_{m\in\mathcal L}
  \sum_{p\in\mathcal P_{o,m,\ell,d}}\psi_{m,\ell}^{p}
 +\sum_{m\ne\ell}\sum_{r\in\mathcal R_{m,\ell}}\varphi_{m,\ell}^{r},
 &&\ell\in\mathcal L,\label{eq:path_balance}\\
 &\sum_{(o,d)\in\mathcal D}\sum_{\ell,m\in\mathcal L}
  \sum_{p\in\mathcal P_{o,\ell,m,d}}b_i^p\psi_{\ell,m}^{p}
 \le\sum_{s\in\mathcal S}\mu_{i,s}z_{i,s},
 &&i\in\mathcal H,\label{eq:path_capacity}\\
 &\sum_{(o,d)\in\mathcal D}\sum_{\ell,m\in\mathcal L}
  \sum_{p\in\mathcal P_{o,\ell,m,d}}a_{u,v}^p\psi_{\ell,m}^{p}
 +\sum_{\ell\ne m}\sum_{r\in\mathcal R_{\ell,m}}a_{u,v}^r\varphi_{\ell,m}^{r}
 \le M_{u,v} y_{u,v},
 &&(u,v)\in\mathcal A,\label{eq:path_activation}\\
 &\sum_{\ell\in\mathcal L}I_\ell
 +\sum_{(o,d)\in\mathcal D}\sum_{\ell,m\in\mathcal L}
 \sum_{p\in\mathcal P_{o,\ell,m,d}}\tau_p\psi_{\ell,m}^{p}
 +\sum_{\ell\ne m}\sum_{r\in\mathcal R_{\ell,m}}\tau_r\varphi_{\ell,m}^{r}
 \le\Gamma,\label{eq:path_full_fleet}\\
 &x_v,z_{i,s},y_{u,v},\alpha_{o,\ell,m,d}\in\{0,1\},
 \qquad I_\ell,\Gamma,\psi_{\ell,m}^{p},\varphi_{\ell,m}^{r}\ge0.
 \label{eq:path_domains}
\end{align}
\end{subequations}
All variables in \eqref{eq:path_domains} range over their previously defined index sets, and $\rho$ and $\Gamma$ are continuous.
The direct-courier flow in \eqref{eq:path_objective} is nonnegative by \eqref{eq:path_linearization} and \eqref{eq:path_assign}.

\section{Solution Method}
\label{sec:path_solution}

We decompose \eqref{eq:path_full} into an integer strategic master problem and a continuous conic operating subproblem.
Column generation handles the service and repositioning paths within the subproblem.
This follows the Benders--column-generation architecture of \citet{bayramHubNetwork2023TransportationScience}, with pricing adapted to battery feasibility and courier--drone coordination.
All candidate paths are defined independently of the strategic decisions, which enter the operating constraints only through affine right-hand sides.

\subsection{Master Problem}
\label{subsec:path_master}

Let $\mathcal X$ contain the binary strategic decisions satisfying \eqref{eq:path_install}--\eqref{eq:path_colocation}.
For a fixed design $(x,z,y,\alpha)$, let $Q(x,z,y,\alpha)$ be its minimum fleet and delivery operating cost, including the direct-courier flow.
The master uses a continuous variable $\zeta$ to approximate this value:
\begin{equation}
\begin{aligned}
 \min_{x,z,y,\alpha,\zeta}\quad
 &\sum_{\ell\in\mathcal L}f_\ell^{\mathrm{base}}x_\ell
 +\sum_{i\in\mathcal H}\sum_{s\in\mathcal S}f_{i,s}^{\mathrm{locker}}z_{i,s}
 +\sum_{k\in\mathcal E}f_k^{\mathrm{swap}}x_k
 +\sum_{(u,v)\in\mathcal A}f_{u,v}^{\mathrm{link}}y_{u,v}
 +\zeta\\
 \text{s.t.}\quad
 &(x,z,y,\alpha)\in\mathcal X,\qquad \zeta\ge0,\\
 &\zeta\ge L^q(x,z,y,\alpha),\qquad q\in\mathcal C.
\end{aligned}
\label{eq:path_master}
\end{equation}
Here $\mathcal C$ indexes the Benders cuts generated so far, with $L^q$ defined in Section~\ref{subsec:path_inequalities}.
The initial operating-cost bound in that section can be included before the first iteration.
The facility, capacity-level, link, and assignment variables remain binary; all fleet, availability, and flow decisions are evaluated in the subproblem.

\subsection{Benders Subproblem}
\label{subsec:path_subproblem}

For an integer master solution $(\widehat x,\widehat z,\widehat y,\widehat\alpha)$, the subproblem is
\begin{equation}
\begin{aligned}
 Q(\widehat x,\widehat z,\widehat y,\widehat\alpha)
 =\min_{\rho,I,\Gamma,\psi,\varphi}\quad
 &f^{\mathrm{fleet}}\Gamma
 +\sum_{(o,d)\in\mathcal D}\sum_{\ell,m\in\mathcal L}
  \sum_{p\in\mathcal P_{o,\ell,m,d}}c_p^{\mathrm{service}}\psi_{\ell,m}^{p}\\
 &+\sum_{\ell\ne m}\sum_{r\in\mathcal R_{\ell,m}}
  c_r^{\mathrm{reposition}}\varphi_{\ell,m}^{r}\\
 &+\sum_{(o,d)\in\mathcal D}c_{o,d}^{\mathrm{rider}}
  \left(h_{o,d}-\sum_{\ell,m\in\mathcal L}
  \sum_{p\in\mathcal P_{o,\ell,m,d}}\psi_{\ell,m}^{p}\right)\\
 \text{s.t.}\quad
 &\eqref{eq:path_linearization},\ \eqref{eq:path_fleet},
  \eqref{eq:path_rho},\ \eqref{eq:path_full_cone},
  \eqref{eq:path_balance}--\eqref{eq:path_activation},\\
 &(x,z,y,\alpha)=(\widehat x,\widehat z,\widehat y,\widehat\alpha),\\
 &I,\Gamma,\psi,\varphi\ge0.
\end{aligned}
\label{eq:path_subproblem}
\end{equation}
The zero-flow solution $\rho=I=\Gamma=\psi=\varphi=0$ serves all orders by couriers and is feasible for every strategic design.
Hence the decomposition requires optimality cuts only.

The conic transformation has already been incorporated in \eqref{eq:path_full_cone}.
The subproblem has an attained optimum and an attained dual optimum, including designs that force some availabilities to zero.
To see this, eliminate $\Gamma$ and $I$ at their minimum-cost values.
The resulting objective is a convex differentiable function of $\rho$ plus linear flow costs over a polyhedron, with queueing term $f^{\mathrm{fleet}}\sum_\ell\rho_\ell/(1-\rho_\ell)$.
The feasible all-courier cost bounds the useful fleet and keeps an optimum away from $\rho_\ell=1$.
Polyhedral first-order optimality conditions then provide finite multipliers for the conic formulation.

For pricing and cut generation, use the nonnegative multipliers in Table~\ref{tab:path_duals}.
Let $\pi_\ell$ be the unrestricted multiplier on base outflow minus inflow in \eqref{eq:path_balance}.
The multiplier on \eqref{eq:path_fleet} can be chosen as $f^{\mathrm{fleet}}$, since $\Gamma$ is continuous and appears only in its positive linear cost and that constraint.

\begin{table}[htbp]
\centering
\caption{Nonnegative operating-subproblem multipliers.}
\label{tab:path_duals}
\begin{tabular}{ll}
\toprule
Constraint & Multiplier\\
\midrule
Pair-flow upper bound \eqref{eq:path_pair} & $\beta_{o,\ell,m,d}$\\
Source-flow upper bound \eqref{eq:path_source_upper} & $\gamma_{o,\ell,d}$\\
Source-flow lower bound, oriented as \eqref{eq:path_source_lower} & $\delta_{o,\ell,d}$\\
Availability bound $\rho_\ell\le x_\ell$ & $\eta_\ell$\\
Locker capacity \eqref{eq:path_capacity} & $\nu_i$\\
Link activation \eqref{eq:path_activation} & $\omega_{u,v}$\\
\bottomrule
\end{tabular}
\end{table}

\subsection{Column Generation}
\label{subsec:path_cg}

The restricted subproblem replaces the complete path sets by subsets $\widehat{\mathcal P}_{o,\ell,m,d}$ and $\widehat{\mathcal R}_{\ell,m}$.
It remains feasible even with an empty path pool because direct courier delivery is available for every order.
Solving it yields feasible operations and dual prices for the current design.

The reduced cost of a service column is
\begin{align}
 \overline c_{\ell,m}^{p}
 ={}&c_p^{\mathrm{service}}-c_{o,d}^{\mathrm{rider}}
 +f^{\mathrm{fleet}}\tau_p+\pi_\ell-\pi_m
 +\beta_{o,\ell,m,d}+\gamma_{o,\ell,d}-\delta_{o,\ell,d}\notag\\
 &+\sum_{i\in\mathcal H}\nu_i b_i^p
 +\sum_{(u,v)\in\mathcal A}\omega_{u,v}a_{u,v}^p,
 \qquad p\in\mathcal P_{o,\ell,m,d}.
 \label{eq:path_rc_service}
\end{align}
The term $-c_{o,d}^{\mathrm{rider}}$ follows from the reduction in direct-courier flow when one additional order uses a service column.
The reduced cost of an empty column is
\begin{equation}
 \overline c_{\ell,m}^{r}
 =c_r^{\mathrm{reposition}}+f^{\mathrm{fleet}}\tau_r
 +\pi_\ell-\pi_m
 +\sum_{(u,v)\in\mathcal A}\omega_{u,v}a_{u,v}^{r}.
 \label{eq:path_rc_empty}
\end{equation}
The pricing problems are
\begin{equation}
 \min_{p\in\mathcal P_{o,\ell,m,d}}\overline c_{\ell,m}^{p}
 \quad\text{for each }(o,d),\ell,m,
 \qquad
 \min_{r\in\mathcal R_{\ell,m}}\overline c_{\ell,m}^{r}
 \quad\text{for each }\ell\ne m.
 \label{eq:path_pricing}
\end{equation}
Negative-reduced-cost paths are added and the restricted subproblem is reoptimized.
Pricing examines the complete candidate graph, including currently closed facilities and links, whose use is controlled by the activation constraints.
Once every pricing value is nonnegative, the restricted dual is feasible for all columns and the operating subproblem is solved.

For service pricing, enumerate the pickup and delivery lockers $i\ne j$ and optimize the energy-feasible movements between the two bases and these locker actions.
Before pickup, a label records elapsed time, accumulated flight and link-price cost, and remaining flight-time budget.
The budget resets at replenishment nodes and is preserved at lockers.
At a replenishment node all labels have a full battery, so a label with no greater elapsed time and no greater accumulated cost dominates another label there.
At a locker, dominance additionally requires no smaller remaining budget.
When pickup occurs after a drone prefix of duration $t$, charge $f^{\mathrm{fleet}}\mathbb E[\max\{t,T_{o,i}\}]$, the first-leg courier cost, and the locker price $\nu_i$.
The expectation is nondecreasing in $t$, so this dominance rule preserves an optimal continuation.

After pickup, all duration costs are additive: each flight contributes $(c^{\mathrm{drone}}+f^{\mathrm{fleet}})t_{u,v}^{\mathrm{fly}}+\omega_{u,v}$, and each replacement contributes $f^{\mathrm{fleet}}t_k^{\mathrm{sw}}$.
The delivery action adds the locker price $\nu_j$ and expected terminal courier cost $(1-\sigma_{d,j})c_{j,d}^{\mathrm{rider}}$.
Given handling times enter the same additive duration cost.
The final replacement at $m$ is included.
Empty pricing applies these additive movement and replacement prices to routes containing only replenishment nodes.

The labeling search may allow walks and remove repeated-node cycles from its output.
All movement and replenishment increments are nonnegative, and a completed replenishment restores the same battery state.
Thus a within-stage cycle cannot improve a label or its best completion, and an optimal priced walk has a stagewise cycle-free representative in the specified path set.
Identical labels are merged, including those induced by zero-time connections.
This implements the dominance rule without imposing additional visited-set restrictions.

\subsection{Valid Inequalities}
\label{subsec:path_inequalities}

At iteration $q$, let $(\widehat x,\widehat z,\widehat y,\widehat\alpha)$ be the tested design and $\widehat Q$ its fully evaluated operating cost.
Use optimal multipliers obtained after all pricing problems have been checked.
The Benders cut is
\begin{align}
 \zeta\ge L^q(x,z,y,\alpha)
 :={}&\widehat Q
 -\sum_{\ell\in\mathcal L}\widehat\eta_\ell(x_\ell-\widehat x_\ell)
 -\sum_{i\in\mathcal H}\sum_{s\in\mathcal S}
 \mu_{i,s}\widehat\nu_i(z_{i,s}-\widehat z_{i,s})\notag\\
 &-\sum_{(u,v)\in\mathcal A}M_{u,v}\widehat\omega_{u,v}(y_{u,v}-\widehat y_{u,v})\notag\\
 &+\sum_{(o,d)\in\mathcal D}\sum_{\ell,m\in\mathcal L}
 h_{o,d}(\widehat\delta_{o,\ell,d}-\widehat\beta_{o,\ell,m,d})
 (\alpha_{o,\ell,m,d}-\widehat\alpha_{o,\ell,m,d}).
 \label{eq:path_benders_cut}
\end{align}

\begin{proposition}
Constraint \eqref{eq:path_benders_cut} is valid for every strategic design in $\mathcal X$ and is tight at its generating design.
\end{proposition}
\begin{proof}
The complete-column dual feasible region is independent of the strategic design.
Keeping its optimal multipliers fixed yields an affine lower bound on the operating value by weak duality.
The right-hand-side derivatives are $-\eta_\ell$ for $x_\ell$, $-\mu_{i,s}\nu_i$ for $z_{i,s}$, $-M_{u,v}\omega_{u,v}$ for $y_{u,v}$, and $h_{o,d}(\delta_{o,\ell,d}-\beta_{o,\ell,m,d})$ for $\alpha_{o,\ell,m,d}$.
Strong duality makes that lower bound equal to $\widehat Q$ at the generating design, giving \eqref{eq:path_benders_cut}.
\end{proof}

An initial master inequality follows from a lower bound on each assigned pair's per-order cost.
Define
\begin{equation}
 \underline c_{o,\ell,m,d}
 =\min_{p\in\mathcal P_{o,\ell,m,d}}
 \{c_p^{\mathrm{service}}+f^{\mathrm{fleet}}\tau_p\},
 \qquad
 \Delta_{o,\ell,m,d}
 =h_{o,d}\max\{0,c_{o,d}^{\mathrm{rider}}-\underline c_{o,\ell,m,d}\}.
 \label{eq:path_saving_bound}
\end{equation}
For an empty path set, take $\underline c_{o,\ell,m,d}=+\infty$ and $\Delta_{o,\ell,m,d}=0$.
Since an assigned OD can send at most $h_{o,d}$ orders through its base pair, while idle inventory and repositioning have nonnegative costs,
\begin{equation}
 \zeta\ge
 \sum_{(o,d)\in\mathcal D}h_{o,d}c_{o,d}^{\mathrm{rider}}
 -\sum_{(o,d)\in\mathcal D}\sum_{\ell,m\in\mathcal L}
 \Delta_{o,\ell,m,d}\alpha_{o,\ell,m,d}
 \label{eq:path_initial_cut}
\end{equation}
is a valid initial lower bound.
The coefficients can be computed by the service pricing procedure with zero locker and link prices.
Any valid lower estimate of $\underline c_{o,\ell,m,d}$ also gives a valid, possibly weaker, inequality.

The activation coefficient in \eqref{eq:path_activation} is justified separately for every fixed design.
A service route traverses a given link at most once in each of its three stages, and total service flow is at most $\sum_{(o,d)}h_{o,d}$.
Remove directed cycles from the base-level empty flow while keeping service flows unchanged.
The remaining empty flow decomposes into imbalance paths of at most $|\mathcal L|-1$ base-to-base movements, and their total source imbalance is at most total service flow.
Each empty route traverses any given physical link at most once.
The combined link usage therefore has an optimal representative bounded by $(|\mathcal L|+2)\sum_{(o,d)}h_{o,d}$, as used in the model.
Cycle removal preserves all base balances and weakly decreases capacity use and fleet cost.

For a link incident to locker $i$, an additional activation inequality is
\begin{equation}
 \sum_{(o,d)\in\mathcal D}\sum_{\ell,m\in\mathcal L}
 \sum_{p\in\mathcal P_{o,\ell,m,d}}a_{u,v}^p\psi_{\ell,m}^{p}
 +\sum_{\ell\ne m}\sum_{r\in\mathcal R_{\ell,m}}a_{u,v}^r\varphi_{\ell,m}^{r}
 \le \left(\max_{s\in\mathcal S}\mu_{i,s}\right)y_{u,v}.
 \label{eq:path_locker_link_cut}
\end{equation}
Every traversal of that link requires the designated handling event at $i$.
For a link joining two lockers, the inequality applies to both endpoints.
Taking the minimum of these bounds and $M^0$ gives \eqref{eq:path_link_bound}.
The same coefficients $M_{u,v}$ are used in the full model, operating subproblem, and Benders cuts.

The algorithm alternates between the integer master and complete operating evaluations by column generation.
The master objective provides a lower bound, and each feasible design with its operating solution provides an upper bound.
A Benders cut is added after complete pricing certifies the subproblem value.
The algorithm terminates when the upper and lower bounds meet the prescribed tolerance.
With exact subproblem solutions and exact pricing, the finite strategic design set and cuts that are tight at their generating designs give finite convergence.

\label{sec:solution_methods}
The model separates naturally into integer strategic decisions and continuous queueing and routing decisions.
Our decomposition combines a path-based conic operating subproblem with a Benders master, following the conic-duality and column-generation architecture of \citet{Bayram2023}.
Here the recourse coefficients incorporate pickup synchronization and a closed-fleet planning relationship, and the design includes a shared fixed charge for every directed physical flight link.
The linear lift in Proposition~\ref{prop:lift} provides the affine design dependence required by the decomposition.

\subsubsection{Master problem}

Let $Q(x,n,y,B)$ be the minimum operational part of \eqref{eq:full-objective}, including $C^0$, for a fixed strategic design.
Given a collection $\mathcal C$ of valid affine recourse lower bounds $L^k$, the master problem is
\begin{equation}
 \begin{aligned}
 \min_{x,n,y,B,\zeta}\quad&C^{\mathrm{fix}}(x,n,y)+\zeta\\
 \text{s.t.}\quad&(x,n,y,B)\in\mathcal X,\qquad\zeta\ge0,\\
 &\zeta\ge L^k(x,n,y,B),\qquad k\in\mathcal C.
 \end{aligned}
 \label{eq:master}
\end{equation}
Only the cost epigraph $\zeta$ is continuous in the master.
The remaining variables are the facility, cabinet-count, flight link, and OD-assignment decisions.
The bound $\zeta\ge0$ follows from nonnegative actual fulfillment and fleet costs.
All queueing and path-flow variables remain in the operating subproblem.

\subsubsection{Subproblem}

For an integer master solution $(\bar x,\bar n,\bar y,\bar B)$, substitute those values into \eqref{eq:full-cone}--\eqref{eq:full-nonnegative} and minimize the operational cost.
This is a second-order cone program.
For an assigned OD--base pair, it is computationally convenient to use the reduced equality $\Theta_{od,\ell m}=\lambda_{od}\rho_\ell$ and to eliminate inactive groups and their source slacks.
Proposition~\ref{prop:lift} ensures that this reduction has the same fixed-design value.
The full lifted form is retained for deriving globally valid cuts.

The restricted subproblem includes subsets $\widehat{\mathcal P}_{od}^{\ell m}\subseteq\mathcal P_{od}^{\ell m}$ and $\widehat{\mathcal R}_{\ell m}\subseteq\mathcal R_{\ell m}$.
Its solution supplies feasible operations for the current design and therefore an upper bound on $Q$.
The zero-flow all-courier solution makes every restricted subproblem feasible, including one with an initially empty path pool.
Initial energy-feasible paths can be added to obtain informative starting solutions.

\subsubsection{Pricing subproblem}

Associate free multipliers $\alpha_{od,\ell}$ with the source equations in the orientation $\lambda_{od}\rho_\ell-\sum_m\Theta_{od,\ell m}-w_{od,\ell}=0$.
Associate free multipliers $\pi_\ell$ with base outflow minus inflow.
Let $u_h,\eta_{uv},\nu_\ell,\sigma_\ell\ge0$ be the multipliers for cabinet capacity, flight link activation, empty-clock bounds, and $\rho_\ell\le\bar\rho x_\ell$, respectively.
Finally, let $\gamma_{od,\ell m}\ge0$ and $\beta_{od,\ell}\ge0$ be the multipliers for \eqref{eq:group-lift} and \eqref{eq:slack-lift}.
Define
\begin{equation}
 A_\ell^*=\sum_{(o,d)}\lambda_{od}\alpha_{od,\ell}
          -\bar\kappa\nu_\ell+\sigma_\ell,
 \qquad
 h(A)=\min_{0\le r\le\bar\rho}
       \left\{c^D\frac{r}{1-r}+Ar\right\}.
 \label{eq:idle-conjugate}
\end{equation}
The minimizer is
\begin{equation}
 r^*(A)=
 \begin{cases}
 0,&A\ge-c^D,\\
 \min\{\bar\rho,\,1-\sqrt{c^D/(-A)}\},&A<-c^D,
 \end{cases}
 \qquad
 h(A)=c^D\frac{r^*(A)}{1-r^*(A)}+Ar^*(A).
 \label{eq:idle-conjugate-closed}
\end{equation}
This scalar minimization is the exact elimination of the idle cone in the operating dual.

For any multiplier vector satisfying
\begin{align}
 \beta_{od,\ell}&\ge\alpha_{od,\ell},\label{eq:dual-w}\\
 c_p+u_{i(p)}+u_{j(p)}+\sum_{(u,v)}\eta_{uv}a_{uv}^p
 +\pi_\ell-\pi_m-\alpha_{od,\ell}+\gamma_{od,\ell m}&\ge0,
 &&p\in\mathcal P_{od}^{\ell m},\label{eq:dual-service}\\
 c_r+\sum_{(u,v)}\eta_{uv}a_{uv}^r+\pi_\ell-\pi_m+\nu_\ell&\ge0,
 &&r\in\mathcal R_{\ell m},\label{eq:dual-empty}
\end{align}
the dual objective is the affine design function
\begin{align}
 L(x,n,y,B)={}&C^0+\sum_\ell h(A_\ell^*)
 -\sum_{(o,d),\ell}M_{od}\beta_{od,\ell}\notag\\
 &+\sum_{(o,d),\ell,m}M_{od}
  (\beta_{od,\ell}-\gamma_{od,\ell m})B_{od,\ell m}
 -\bar\rho\sum_\ell\sigma_\ell x_\ell\notag\\
 &-\sum_h\mu_hu_hn_h
 -\sum_{(u,v)}U_{uv}\eta_{uv}y_{uv}.
 \label{eq:global-cut}
\end{align}
Maximizing this expression over the dual-feasible multipliers gives $Q$ under strong duality.



For a service group, the pricing problem computes
\begin{equation}
 q_{od,\ell m}(u,\eta)=
 \min_{p\in\mathcal P_{od}^{\ell m}}
 \left\{c_p+u_{i(p)}+u_{j(p)}+
       \sum_{(u,v)}\eta_{uv}a_{uv}^p\right\}.
 \label{eq:service-pricing}
\end{equation}
Its minimum reduced cost is
\begin{equation}
 \bar c_{od,\ell m}=q_{od,\ell m}(u,\eta)
 +\pi_\ell-\pi_m-\alpha_{od,\ell}+\gamma_{od,\ell m}.
 \label{eq:service-reduced-cost}
\end{equation}
An empty-path pricing problem computes
\begin{equation}
 q_{\ell m}^{E}(\eta)=\min_{r\in\mathcal R_{\ell m}}
 \left\{c_r+\sum_{(u,v)}\eta_{uv}a_{uv}^r\right\},
 \label{eq:empty-pricing}
\end{equation}
with minimum reduced cost $q_{\ell m}^{E}(\eta)+\pi_\ell-\pi_m+\nu_\ell$.
Every negative-reduced-cost path is eligible to enter the restricted subproblem.
Pricing uses all candidate facilities and physical links, including currently disabled links whose flow is controlled by \eqref{eq:airline-activation}.
This preserves the dual rows for the complete design-independent column universe.

The pricing calculation exploits the complete replacement at every visited node in $\mathcal K$.
Before pickup, a label at such a node $a$ records elapsed time $t$ and accumulated flight-distance and link-price cost $z$; all labels at $a$ have the same full battery resource.
The initial label at $\ell$ is $(0,0)$.
A station-to-station movement adds its flight and replacement duration to $t$, and its flight-distance charge and $\eta$ price to $z$.
A label with no greater $t$ and no greater $z$ dominates another label at the same station, because every continuation begins with the same resource and the synchronization transform is nondecreasing in arrival time.
Let $\mathcal F_{\ell a}$ denote the resulting nondominated prefix labels.
Nonnegative increments ensure that cycles cannot produce a new nondominated label; Proposition~\ref{prop:finite-paths} therefore gives a finite exact calculation.

After pickup, all time costs are additive.
Station-to-station movements have price $c^{\mathrm{fly}}d_{uv}+c^D(t_{uv}+s_v)+\eta_{uv}$, where $s_v$ is the replacement duration at the arrival station.
Their shortest-path closure accounts for any number of intermediate replacements.
For fixed last pre-pickup station $a$, pickup locker $i$, delivery locker $j$, and return base $m$, let $D_{aijm}(\eta)$ be the minimum subsequent flight, fleet-time, and link-price cost, starting at $i$ with resource $\bar e-d_{ai}$.
It is obtained by comparing a direct loaded link $i\to j$ with a route that first visits a replacement station, and then combining the station closure with delivery and return movements.
Enumerating the stations adjacent to the two locker actions enforces the residual-energy sums between successive replacements.
The terminal replacement at $m$ is included in $D_{aijm}$.

The service pricing value can consequently be evaluated as
\begin{equation}
\begin{aligned}
 q_{od,\ell m}(u,\eta)=
 \min_{\substack{i,j\in\mathcal H,\ i\ne j;\ a\in\mathcal K\\
                 (a,i)\in\mathcal A;\ (t,z)\in\mathcal F_{\ell a}}}
 \bigl\{&z+c^{\mathrm{fly}}d_{ai}+\eta_{ai}
       +c^D\mathbb E[\max\{t+t_{ai},T_{oi}\}]\\
       &+D_{aijm}(\eta)+c_{oi}^F+\bar c_{dj}-c_{od}^R+u_i+u_j\bigr\}.
\end{aligned}
\label{eq:pricing-factorization}
\end{equation}
Energy-infeasible combinations have $D_{aijm}=+\infty$.
This factorization shares post-pickup shortest-path calculations across OD groups while retaining the two-dimensional pre-pickup time--cost tradeoff.
Empty pricing uses the same nonnegative-price station closure from $\ell$ to $m$, including replacement at every arrival.

\subsection{Valid Inequalities}
\label{sec:valid_inequalities_integrated}


After complete pricing, source-slack and group multipliers can be completed explicitly as
\begin{align}
 \beta_{od,\ell}&=\max\{\alpha_{od,\ell},0\},\label{eq:beta-completion}\\
 \gamma_{od,\ell m}&=
 \max\{0,\alpha_{od,\ell}-\pi_\ell+\pi_m-q_{od,\ell m}(u,\eta)\}.
 \label{eq:gamma-completion}
\end{align}
For an empty service-path set, its minimum cost is $+\infty$ and the corresponding completed $\gamma$ is zero.
The empty reduced-cost conditions are checked over the full replacement network.
To make this certificate explicit, first replace each empty-clock multiplier by
\begin{equation}
 \nu_\ell\leftarrow\nu_\ell+
 \max\left\{0,\ \max_{m\ne\ell:\,\mathcal R_{\ell m}\ne\varnothing}
 [\pi_m-\pi_\ell-\nu_\ell-q_{\ell m}^E(\eta)]\right\}.
 \label{eq:empty-dual-repair}
\end{equation}
The increment is zero when there is no reachable return base or every empty reduced cost is already nonnegative.
If the fixed-design solver eliminated inactive groups, let $a_{od,\ell m}$ be the multiplier on an assigned group in the reduced dual orientation $a_{od,\ell m}\le q_{od,\ell m}(u,\eta)$.
Embed that multiplier by setting
\begin{equation}
 \alpha_{od,\ell}=
 \min\{a_{od,\ell m},q_{od,\ell m}(u,\eta)\}+\pi_\ell-\pi_m
 \quad\text{for the selected pair $(\ell,m)$},
 \label{eq:alpha-embedding}
\end{equation}
and set $\alpha_{od,\ell}=0$ at unselected sources.
Equations~\eqref{eq:beta-completion}--\eqref{eq:gamma-completion} then complete the source-slack and group prices for every candidate pair.
Finally, recompute $A_\ell^*$ and $h(A_\ell^*)$ using the repaired $\nu_\ell$ before evaluating the cut.

\begin{proposition}[Global Benders validity]
\label{prop:cut-validity}
Suppose \eqref{eq:dual-service} and \eqref{eq:dual-empty} hold over the complete candidate path sets, with the stated multiplier signs and \eqref{eq:dual-w}.
Then $\zeta\ge L(x,n,y,B)$ is a valid optimality cut for every strategic design in $\mathcal X$.
If the multipliers are optimal for the full operating subproblem at a design, the cut equals its recourse value there.
\end{proposition}
\begin{proof}
Form the Lagrangian of the lifted operational problem using the stated orientations.
The coefficient of each service or empty flow is its reduced cost in \eqref{eq:dual-service} or \eqref{eq:dual-empty}, and the coefficient of $w_{od,\ell}$ is $\beta_{od,\ell}-\alpha_{od,\ell}$.
Their nonnegativity makes the infimum over these nonnegative variables equal to zero.
Minimizing over the idle-cone variables gives $h(A_\ell^*)$.
The remaining right-hand-side terms are exactly \eqref{eq:global-cut}.
Weak duality therefore establishes $L\le Q$ for every design, since the dual feasibility region contains no strategic variables.
Strong duality gives equality at an optimal multiplier vector.
\end{proof}

The algorithm solves the integer master, evaluates a feasible operating solution for its candidate design, and solves the restricted operational problem with repeated pricing.
A full-network dual certificate generates \eqref{eq:global-cut}, which is added to the master before the next design is selected.
The best verified complete solution provides an upper bound, while the master provides a global lower bound.
The algorithm terminates when their prescribed absolute or relative gap is met.
With exact subproblem optimization and exact pricing, the finite strategic design set and supporting cuts provide finite convergence.
Under a computation limit, the algorithm retains its best feasible solution and its valid global lower bound.
Any numerical reduced-cost violation is accounted for in the certified lower bound before the associated cut is used.

The decomposition therefore preserves an integer design master, continuous availability decisions, and detailed energy-feasible path allocation.
Locker capacities and fixed flight link decisions couple different OD missions through shared physical resources, while the complete-column dual supplies the common basis for pricing and Benders cuts.


\color{black}

\begingroup






The variables $\chi_{odijr}$ determine how much OD demand is offered to each corridor and availability level, while constraint~\eqref{con:od_reliability_match} connects this allocation to $w$ and constraints~\eqref{con:source_flow}--\eqref{con:loaded_flow} generate realized drone flow.
Consequently, a network cut that forces capacity for all exogenous market demand is not valid for the current model without an additional minimum drone-service requirement.
We nevertheless retain the max-flow/min-cut construction because it may become useful when the final experimental specification imposes a minimum drone-service level.

Suppose a future formulation requires at least $D^{\mathrm{req}}(\boldsymbol x)$ units of drone flow and let $G^F(\boldsymbol x,\boldsymbol y,\boldsymbol z)$ be an auxiliary flow graph whose node and arc capacities are induced by the selected facilities, smart-locker capacities, and physical arcs.
For any cut set $\mathcal K$ that separates the required sources from the sinks, the corresponding candidate inequality is
\begin{align}
\sum_{e\in\mathcal K}
\operatorname{Cap}_{e}(\boldsymbol x,\boldsymbol y,\boldsymbol z)
\ge
D^{\mathrm{req}}(\boldsymbol x).
\label{con:candidate_min_cut}
\end{align}
If an integer network design cannot route the required flow in $G^F$, a minimum-cut calculation identifies a violated inequality of the form~\eqref{con:candidate_min_cut}.
Such cuts could be separated at integer incumbents in a branch-and-cut implementation.

Constraint~\eqref{con:candidate_min_cut} is retained as a candidate strengthening idea, not claimed as a valid inequality for the current partial-service formulation.
Its final right-hand side and capacity definitions must be derived only after a positive service lower bound and any separate swapping-station throughput capacities used in the auxiliary cut graph are specified.

% The max-flow/min-cut idea may also be used as a fast screening test before solving the full conic model.
% The aggregate energy-sufficiency constraints do not by themselves define station-level swapping-throughput capacities for the auxiliary cut graph.

\section{Computational Study and Managerial Analysis}
\label{sec:experiments}

\subsection{Data and Instance Construction}
\label{subsec:data_instances}

The baseline instances are constructed from an anonymized 2023 sample of 592,667 delivery orders in Nanshan District, Shenzhen, supplied by Meituan for the 2026 City Low-Altitude Aviation Network business-analysis competition \citep{Meituan2026NanshanDataset}.
The source replaces calendar dates with seven sequential day identifiers and applies anonymization and grid processing to order and coordinate fields; no customer identifiers are used.
For the baseline 16-km-by-16-km study window, we retain orders whose pickup and delivery endpoints both lie within the window and use the 12:00--13:00 records to construct demand.
We aggregate pickup and delivery endpoints to 0.5-km grid cells, define each nonzero cell pair as an OD market, and divide its seven-day order count by seven one-hour observation periods to obtain $h_{od}$.
This procedure retains 56,045 records, creates 11,004 nonzero OD markets, and yields total demand of 8,006.43 orders per hour.
Candidate locations are generated from 08:00--24:00 activity on a 0.25-km grid: joint pickup--drop-off activity for smart lockers, pickup activity for drone bases, and pickup--drop-off midpoint activity for battery-swapping stations.
The baseline candidate set contains 34 smart lockers, 5 drone bases, and 8 battery-swapping stations; candidate-density experiments regenerate this set from the same records and preprocessing rules.
The order data identify spatial demand and candidate-site activity only.
The MNL coefficients in~\eqref{eq:choice_utilities}, the 30-minute home-delivery attribute, 20-minute locker-delivery attribute, locker-price discount, maximum walking distance, delivery-cost coefficients, availability levels, energy parameters, service rates, flight speed, infrastructure costs, and other operating parameters are transparent scenario inputs rather than empirical estimates from these records.
All sensitivity results therefore distinguish data-derived demand and candidate locations from assumed behavioral, cost, and operating parameters.

\subsection{Solution Quality and Cost Reporting}

For every instance and service scope, the computational study reports the solver status, incumbent objective value, best bound, relative optimality gap, and runtime.
The objective value is reported as modeled platform fulfillment cost rather than revenue or profit.
Its components are facility, smart-locker capacity, physical-connection, route-operation, fleet, and courier-fulfillment costs, using the same monetary-per-time unit as objective~\eqref{obj:integrated}.
The all-courier counterfactual is $\sum_{o,d\in\mathcal N}f_{od}^{\mathrm{rider}}h_{od}$.
Modeled cost saving is the all-courier counterfactual minus the optimized objective value and is not interpreted as profit.

\subsection{Service Accounting}

The following quantities are computed directly from the variables in MLND:
\begin{subequations}
\label{eq:reported_service_components}
\begin{align}
s^{\mathrm{dr}}
&=
\sum_{o,d\in\mathcal N}
\sum_{(i,j)\in\mathcal A^{HH}}
\sum_{r\in\mathcal R}
\rho_r\chi_{odijr}
=
\sum_{(i,j)\in\mathcal A^{HH}}\psi_{ij},
\label{eq:reported_drone_service}
\\
s^{\mathrm{lk}}
&=
\sum_{o,d\in\mathcal N}
\sum_{j\in\mathcal H_d}
\eta_{odj}\epsilon_{odj},
\label{eq:reported_locker_service}
\\
s^{\mathrm{hm}}
&=
s^{\mathrm{dr}}-s^{\mathrm{lk}},
\label{eq:reported_assisted_home}
\\
s^{\mathrm{fb}}
&=
\sum_{o,d\in\mathcal N}h_{od}-s^{\mathrm{dr}}.
\label{eq:reported_fallback}
\end{align}
\end{subequations}
The first quantity is the rate realized through loaded drone corridors.
The second is the successful locker-pickup rate after applying customer choice at the displayed destination locker.
The third is the successful drone-assisted home-delivery rate.
The fourth is the rate delivered directly by couriers after nonassignment or unsuccessful drone realization.
These quantities must satisfy $s^{\mathrm{lk}}+s^{\mathrm{hm}}+s^{\mathrm{fb}}=\sum_{o,d}h_{od}$ in every reported solution.
For the home-delivery-assistance restriction, $s^{\mathrm{lk}}=0$ and $s^{\mathrm{hm}}=s^{\mathrm{dr}}$.
For the locker-only restriction, $s^{\mathrm{lk}}=s^{\mathrm{dr}}$ and $s^{\mathrm{hm}}=0$.

\subsection{Network-Use Metrics}

When total physical flow is positive, the share routed through a battery-swapping station is
\begin{align}
r^{\mathrm{swap}}
=
\frac{
\displaystyle
\sum_{(u,v)\in\mathcal A^{F}}
\sum_{k\in\mathcal E_{uv}}
\varphi_{uv}^{k}
}{
\displaystyle
\sum_{(u,v)\in\mathcal A^{F}}
\psi_{uv}
}.
\label{eq:swap_flow_ratio}
\end{align}
The numerator counts physical flow assigned to one-swap route options, and the denominator counts all physical drone flow.
When total loaded flow is positive, the chaining share is
\begin{align}
r^{\mathrm{chain}}
=
\frac{
\displaystyle
\sum_{i\in\mathcal H}\psi_{s_i i}
}{
\displaystyle
\sum_{(i,j)\in\mathcal A^{HH}}\psi_{ij}
}.
\label{eq:chaining_ratio}
\end{align}
Its numerator counts realized orders dispatched from virtual idle states, while its denominator is total realized loaded service.
Either ratio is reported as not applicable when its denominator is zero.
Facility counts, selected smart-locker capacities, activated physical connections, fleet size $\Gamma$, and the three fleet-population components in constraint~\eqref{con:fleet_size} are reported separately.

\subsection{Formulation-Matched Comparisons}

The mixed-delivery MLND, home-delivery-assistance restriction, and locker-only restriction are compared under the same OD demand, candidate infrastructure, cost coefficients, travel times, and reliability grid.
Each scope is evaluated against the same all-courier counterfactual.
The comparison reports total modeled cost, its components, the four service quantities in~\eqref{eq:reported_service_components}, network topology, and fleet composition.
This common reporting basis isolates the consequences of service scope without introducing a revenue term or reconstructing OD completion from aggregate corridor flow.

\subsection{Sensitivity Analysis}

Sensitivity experiments vary the locker alternative-specific constant $\beta_0^{\mathrm{lk}}$, time, price, and walking sensitivities, quoted delivery times, locker-price discount, maximum walking distance $d^{\mathrm{walk}}$, smart-locker handling-time limit $\tau^H$, virtual-state waiting-time limit $\tau^S$, reliability levels $\rho_r$, usable battery duration $(1-\varepsilon)\Omega/\omega$, infrastructure and fleet cost coefficients, candidate-facility density, and demand scale.
For each change, the study reports both absolute values and changes relative to a clearly identified baseline.
Managerial implications are drawn only from formulation-matched feasible solutions and are stated in terms of modeled cost, realized service, congestion, network structure, and fleet use.
\section{Conclusion}
\label{sec:conclusion}


\bibliographystyle{elsarticle-harv}
\bibliography{MLND,MLND-sources}

\appendix
\begingroup
\section{Complete MLND Formulations}
\label{app:full_model}

This appendix first reports the self-contained mixed-delivery MLND formulation corresponding to Section~\ref{sec:formulation}.
Sections~\ref{app:hd_model} and~\ref{app:lo_model} then specify its home-delivery-assistance and locker-only restrictions, respectively.

\subsection{Sets and Preprocessed Relations}

The physical facilities are drone bases $\mathcal L$, smart lockers $\mathcal H$, and battery-swapping stations $\mathcal E$.
Each smart locker $i\in\mathcal H$ has one virtual idle-drone state $s_i$, so $\mathcal S=\{s_i:i\in\mathcal H\}$, while $\mathcal N$ contains all demand points.
The principal network sets are
\begin{align}
\widetilde{\mathcal L}
={}&
\mathcal L\cup\mathcal S,
\label{app:set_idle}
\\
\mathcal A
={}&
\mathcal A^{LH}
\cup\mathcal A^{SH}
\cup\mathcal A^{HH}
\cup\mathcal A^{HL}
\cup\mathcal A^{HS},
\label{app:set_balance_arcs}
\\
\mathcal A^{F}
={}&
\mathcal A^{LH}
\cup\mathcal A^{HH}
\cup\mathcal A^{HL}.
\label{app:set_physical_arcs}
\end{align}
For each $(u,v)\in\mathcal A^{F}$, the feasible route-option set is
\begin{align}
\widetilde{\mathcal E}_{uv}
=
\begin{cases}
\mathcal E_{uv}\cup\{0\},
& \text{if direct flight from $u$ to $v$ is feasible},\\
\mathcal E_{uv},
& \text{otherwise},
\end{cases}
\label{app:set_route_options}
\end{align}
where $\mathcal E_{uv}\subseteq\mathcal E$ contains the feasible one-swap stations and option $0$ denotes direct flight.
The openable physical connections and eligible idle states are
\begin{align}
\mathcal A^{Y}
={}&
\mathcal A^{F}
\cup
\left\{
(u,k),(k,v):
(u,v)\in\mathcal A^{F},\
k\in\mathcal E_{uv}
\right\},
\label{app:set_openable_arcs}
\\
\widetilde{\mathcal L}_{i}
={}&
\left\{
l\in\mathcal L:(l,i)\in\mathcal A^{LH}
\right\}
\cup\{s_i\},
&&\forall i\in\mathcal H.
\label{app:set_eligible_idle}
\end{align}
The logical arc sets are $\mathcal A^{SH}=\{(s_i,i):i\in\mathcal H\}$ and $\mathcal A^{HS}=\{(i,s_i):i\in\mathcal H\}$.

For each destination $d$, the eligible pickup lockers and the strictly closer eligible lockers are
\begin{align}
\mathcal H_d
={}&
\left\{j\in\mathcal H:d_{dj}\le d^{\mathrm{walk}}\right\},
&&\forall d\in\mathcal N,
\label{app:set_choice_lockers}
\\
\mathcal H_{dj}^{<}
={}&
\left\{j'\in\mathcal H_d:d_{dj'}<d_{dj}\right\},
&&\forall d\in\mathcal N,\ j\in\mathcal H_d.
\label{app:set_closer_lockers}
\end{align}
The normalized systematic utilities are
\begin{subequations}
\label{app:choice_utilities}
\begin{align}
U_{od}^{\mathrm{hm}}
={}&
-\beta^{t}\widehat t_{od}^{\mathrm{hm}},
&&\forall o,d\in\mathcal N,
\label{app:home_utility}
\\
U_{odj}^{\mathrm{lk}}
={}&
\beta_0^{\mathrm{lk}}
-\beta^{t}\widehat t_{odj}^{\mathrm{lk}}
+\beta^{p}\Delta p_{odj}
-\beta^{w}d_{dj},
&&\forall o,d\in\mathcal N,\ j\in\mathcal H_d.
\label{app:locker_utility}
\end{align}
\end{subequations}
The corresponding conditional two-alternative MNL share is
\begin{align}
\eta_{odj}
={}&
\frac{\exp\!\left(U_{odj}^{\mathrm{lk}}\right)}
{\exp\!\left(U_{od}^{\mathrm{hm}}\right)+\exp\!\left(U_{odj}^{\mathrm{lk}}\right)}
=
\frac{1}{1+\exp\!\left[-\left(U_{odj}^{\mathrm{lk}}-U_{od}^{\mathrm{hm}}\right)\right]},
&&\forall o,d\in\mathcal N,\ j\in\mathcal H_d.
\label{app:choice_share}
\end{align}
The utilities and $\eta_{odj}$ are derived before optimization, whereas alternative availability is determined by $x_j$ and $\alpha_{dj}$.
The reliability levels are indexed by $\mathcal R=\{1,\ldots,R\}$ and satisfy $0<\rho_1<\cdots<\rho_R<1$.

\subsection{Parameters and Variables}

The set of smart-locker capacity levels is $\mathcal C$, and the service-rate parameter of level $c\in\mathcal C$ is $\mu_c$.
The exogenous order rate is $h_{od}$ for $o,d\in\mathcal N$, with $h_{od}=0$ for a nonexistent market, $d_{dj}$ is the walking distance from destination $d$ to locker $j$, and $d^{\mathrm{walk}}$ is the maximum admissible walking distance.
The customer-choice inputs are $\beta_0^{\mathrm{lk}}$, $\beta^{t}$, $\beta^{p}$, $\beta^{w}$, $\widehat t_{od}^{\mathrm{hm}}$, $\widehat t_{odj}^{\mathrm{lk}}$, and $\Delta p_{odj}$; the utilities and $\eta_{odj}$ are derived through~\eqref{app:choice_utilities}--\eqref{app:choice_share}.
The maximum average times at smart lockers and virtual holding states are $\tau^H$ and $\tau^S$, respectively.
The travel time of route option $k\in\widetilde{\mathcal E}_{uv}$ is $t_{uv}^{k}$.
The energy parameters are the depletion rate $\omega$, nominal battery energy $\Omega$, and reserve fraction $\varepsilon$; the scalar parameter $\varepsilon$ is distinct from the indexed auxiliary variable $\epsilon_{odj}$.
The constants $M^1_{ij}$, $M^2_{uv}$, $M^3_i$, and $M^4_{odj}$ are used for offered-flow activation, physical-flow activation, smart-locker population activation, and displayed-locker product linearization, respectively.
We set $M^4_{odj}=\rho_Rh_{od}$ for $j\in\mathcal H_d$.
The cost coefficients are $f_u^{F}$ for facilities, $f_{ic}^{C}$ for locker capacity, $f_{uv}^{A}$ for physical connections, $f_{uvk}^{O}$ for route operations, $f^{D}$ for drones, $f_{oi}^{\mathrm{rider}}$, $f_{jd}^{\mathrm{rider}}$, and $f_{od}^{\mathrm{rider}}$ for the corresponding courier movements, and $f_j^{\mathrm{locker}}$ for locker handling.

The binary variables are $x_u$, $y_{uv}$, $z_{ic}$, $\delta_{\ell r}$, and $\alpha_{dj}$.
The nonnegative continuous variables are $\chi_{odijr}$, $w_{\ell i j r}$, $\epsilon_{odj}$, $\psi_{uv}$, $\varphi_{uv}^{k}$, and $\xi_i$.
The nonnegative integer variable $\Gamma$ is the fleet size.
The structural variable $\chi_{odijr}$ retains the market, corridor, and availability-level identity before service realization.
The variable $w_{\ell i j r}$ allocates the same offered pool among eligible idle states.
The variable $\epsilon_{odj}$ is the only product-linearization auxiliary variable and qualifies successful flow for the locker-pickup cost adjustment.

Table~\ref{tab:model_parameters} consolidates all scalar and indexed numerical inputs required to instantiate MLND and its two restrictions; sets, preprocessed relations, and decision variables are excluded.
The third column reports units, admissible domains, or exact preprocessing requirements rather than treating the transparent scenario inputs as empirical estimates.
The home-delivery-assistance restriction omits all customer-choice inputs and derived quantities, $M^4_{odj}$, and $f_j^{\mathrm{locker}}$, whereas the locker-only restriction omits $M^4_{odj}$ and $f_{jd}^{\mathrm{rider}}$; the aggregate energy parameters are shared by all three service scopes.

\begingroup
\small
\renewcommand{\arraystretch}{1.05}
\begin{longtable}{@{}>{\raggedright\arraybackslash}p{0.16\textwidth}>{\raggedright\arraybackslash}p{0.54\textwidth}>{\raggedright\arraybackslash}p{0.22\textwidth}@{}}
\caption{Summary of model parameters}
\label{tab:model_parameters}\\
\toprule
Parameter & Description & Unit or requirement \\
\midrule
\endfirsthead
\multicolumn{3}{c}{Table~\thetable\ continued}\\
\toprule
Parameter & Description & Unit or requirement \\
\midrule
\endhead
\midrule
\multicolumn{3}{r}{Continued on next page}\\
\endfoot
\bottomrule
\endlastfoot
$\mu_c$ & Handling service rate provided by smart-locker capacity level $c$. & Drone handlings per unit time. \\
$\tau^H$ & Maximum admissible average handling time at a smart locker. & Time; positive. \\
$\tau^S$ & Maximum admissible average waiting time in a virtual idle-drone state. & Time; positive. \\
$\omega$ & Mean battery-energy depletion rate per drone while traveling or being handled. & Energy per drone per unit time; positive. \\
$\Omega$ & Nominal energy of a fully charged battery. & Energy per drone; positive. \\
$\varepsilon$ & Required battery-energy reserve fraction. & Dimensionless; $0\le\varepsilon<1$. \\
$f_u^{F}$ & Equivalent facility cost rate for opening physical facility $u\in\mathcal L\cup\mathcal H\cup\mathcal E$. & Currency per unit time. \\
$f_{ic}^{C}$ & Equivalent capacity cost rate for selecting level $c$ at smart locker $i$. & Currency per unit time. \\
$f_{uv}^{A}$ & Equivalent design cost rate for activating physical connection $(u,v)\in\mathcal A^Y$. & Currency per unit time. \\
$f_{uvk}^{O}$ & Route-operation cost per traversal of option $k$ for physical movement $(u,v)$. & Currency per drone movement. \\
$f^{D}$ & Equivalent fleet cost rate per drone. & Currency per drone per unit time. \\
$f_{oi}^{\mathrm{rider}}$ & Courier cost per order from customer origin $o$ to origin locker $i$. & Currency per order. \\
\end{longtable}
\endgroup

\subsection{Objective}

MLND has the following single objective:
\begin{align}
\min\quad
&
\sum_{u\in\mathcal L\cup\mathcal H\cup\mathcal E}f_u^{F}x_u
+
\sum_{i\in\mathcal H}\sum_{c\in\mathcal C}f_{ic}^{C}z_{ic}
+
\sum_{(u,v)\in\mathcal A^{Y}}f_{uv}^{A}y_{uv}
\nonumber\\
&+
\sum_{(u,v)\in\mathcal A^{F}}
\sum_{k\in\widetilde{\mathcal E}_{uv}}
f_{uvk}^{O}\varphi_{uv}^{k}
+f^{D}\Gamma
+
\sum_{o,d\in\mathcal N}f_{od}^{\mathrm{rider}}h_{od}
\nonumber\\
&+
\sum_{o,d\in\mathcal N}
\sum_{(i,j)\in\mathcal A^{HH}}
\sum_{r\in\mathcal R}
\rho_r
\left(
f_{oi}^{\mathrm{rider}}
+f_{jd}^{\mathrm{rider}}
-f_{od}^{\mathrm{rider}}
\right)
\chi_{odijr}
\nonumber\\
&+
\sum_{o,d\in\mathcal N}
\sum_{j\in\mathcal H_d}
\eta_{odj}
\left(
f_j^{\mathrm{locker}}
-f_{jd}^{\mathrm{rider}}
\right)
\epsilon_{odj}.
\label{app:objective}
\end{align}
The first line contains facility, smart-locker capacity, and connection costs.
The second line contains route-operation and fleet costs and then charges every order the direct courier baseline cost.
The third line replaces that baseline for realized drone orders by the corresponding first-mile and last-mile courier costs.
The fourth line replaces the last-mile courier cost by locker handling cost for the fraction that completes locker pickup.

\subsection{Infrastructure and OD Assignment Constraints}

Each open smart locker selects one capacity level, and each activated physical connection requires both endpoints:
\begin{align}
&
\sum_{c\in\mathcal C}z_{ic}
=x_i,
&&\forall i\in\mathcal H,
\label{app:capacity_select}
\\
&
y_{uv}
\le x_u,
&&\forall(u,v)\in\mathcal A^{Y},
\label{app:arc_tail}
\\
&
y_{uv}
\le x_v,
&&\forall(u,v)\in\mathcal A^{Y}.
\label{app:arc_head}
\end{align}
OD demand offered to the drone network satisfies
\begin{align}
&
\sum_{(i,j)\in\mathcal A^{HH}}
\sum_{r\in\mathcal R}
\chi_{odijr}
\le h_{od},
&&\forall o,d\in\mathcal N,
\label{app:od_demand_bound}
\\
&
\sum_{o,d\in\mathcal N}\chi_{odijr}
=
\sum_{\ell\in\widetilde{\mathcal L}_{i}}w_{\ell i j r},
&&\forall(i,j)\in\mathcal A^{HH},\ r\in\mathcal R,
\label{app:od_reliability_match}
\\
&
\sum_{\ell\in\widetilde{\mathcal L}_{i}}
\sum_{r\in\mathcal R}
w_{\ell i j r}
\le M^1_{ij}y_{ij},
&&\forall(i,j)\in\mathcal A^{HH}.
\label{app:corridor_offer_activation}
\end{align}
Constraint~\eqref{app:od_demand_bound} limits the total offered rate of each market.
Constraint~\eqref{app:od_reliability_match} matches the OD allocation and idle-state allocation before service realization for every corridor and availability level.
Constraint~\eqref{app:corridor_offer_activation} activates the total offered pool only when its loaded corridor is selected.

\subsection{Availability, Flows, Routes, and Energy}

The availability selections are
\begin{align}
&
\sum_{r\in\mathcal R}\delta_{l r}
=x_l,
&&\forall l\in\mathcal L,
\label{app:availability_base}
\\
&
\sum_{r\in\mathcal R}\delta_{s_i r}
\le x_i,
&&\forall i\in\mathcal H.
\label{app:availability_virtual}
\end{align}
Offered orders are linked to the selected availability level by
\begin{align}
&
w_{\ell i j r}
\le M^1_{ij}\delta_{\ell r},
&&\forall \ell\in\widetilde{\mathcal L}_{i},\
(i,j)\in\mathcal A^{HH},\
r\in\mathcal R.
\label{app:w_reliability}
\end{align}
The realized source and loaded flows are two marginals of the same $w$ variables:
\begin{subequations}
\label{app:realized_flow_definitions}
\begin{align}
&
\psi_{\ell i}
=
\sum_{j:(i,j)\in\mathcal A^{HH}}
\sum_{r\in\mathcal R}
\rho_r w_{\ell i j r},
&&\forall i\in\mathcal H,\ \ell\in\widetilde{\mathcal L}_{i},
\label{app:source_flow}
\\
&
\psi_{ij}
=
\sum_{\ell\in\widetilde{\mathcal L}_{i}}
\sum_{r\in\mathcal R}
\rho_r w_{\ell i j r},
&&\forall(i,j)\in\mathcal A^{HH}.
\label{app:loaded_flow}
\end{align}
\end{subequations}
The first equality aggregates realized dispatches from each eligible idle state, while the second aggregates realized loaded flow on each corridor.
Together with constraint~\eqref{app:od_reliability_match}, the second equality implies $\psi_{ij}=\sum_{o,d,r}\rho_r\chi_{odijr}$ and therefore preserves the forward realization logic.
Closed-network balance and physical-connection activation require
\begin{align}
&
\sum_{v:(v,u)\in\mathcal A}\psi_{vu}
=
\sum_{v:(u,v)\in\mathcal A}\psi_{uv},
&&\forall u\in\mathcal L\cup\mathcal H\cup\mathcal S,
\label{app:flow_balance}
\\
&
\psi_{uv}
\le M^2_{uv}y_{uv},
&&\forall(u,v)\in\mathcal A^{F}.
\label{app:physical_flow}
\end{align}
Direct and one-swap route flows satisfy
\begin{align}
&
\sum_{k\in\widetilde{\mathcal E}_{uv}}
\varphi_{uv}^{k}
=\psi_{uv},
&&\forall(u,v)\in\mathcal A^{F},
\label{app:route_decomposition}
\\
&
\varphi_{uv}^{k}
\le M^2_{uv}y_{uk},
&&\forall(u,v)\in\mathcal A^{F},\
k\in\mathcal E_{uv},
\label{app:route_swap_in}
\\
&
\varphi_{uv}^{k}
\le M^2_{uv}y_{kv},
&&\forall(u,v)\in\mathcal A^{F},\
k\in\mathcal E_{uv}.
\label{app:route_swap_out}
\end{align}
The route-decomposition equality assigns every physical flow unit to one feasible direct or one-swap option.
The two inequalities activate the access and egress legs whenever a swap-assisted option is used.
Aggregate energy sufficiency requires
\begin{align}
&
\omega
\sum_{i:(l,i)\in\mathcal A^{LH}}
\sum_{k\in\widetilde{\mathcal E}_{li}}
t_{li}^{k}\varphi_{li}^{k}
\le
(1-\varepsilon)\Omega
\sum_{i:(l,i)\in\mathcal A^{LH}}
\sum_{k\in\mathcal E_{li}}
\varphi_{li}^{k},
&&\forall l\in\mathcal L,
\label{app:energy_base}
\\
&
\omega
\left(
\xi_i+
\sum_{v:(i,v)\in\mathcal A^{F}}
\sum_{k\in\widetilde{\mathcal E}_{iv}}
t_{iv}^{k}\varphi_{iv}^{k}
\right)
\le
(1-\varepsilon)\Omega
\sum_{v:(i,v)\in\mathcal A^{F}}
\sum_{k\in\mathcal E_{iv}}
\varphi_{iv}^{k},
&&\forall i\in\mathcal H.
\label{app:energy_locker}
\end{align}
Constraint~\eqref{app:energy_base} bounds the aggregate energy depletion of departures from each physical base by the usable energy supplied through swap-assisted departures.
Constraint~\eqref{app:energy_locker} additionally includes the expected smart-locker handling population $\xi_i$ before applying the same aggregate requirement to departures from locker $i$.
Direct route options contribute to energy consumption but not to the swap count on the right-hand sides.
Standby energy consumption in the logical virtual states is neglected, and individual state-of-charge trajectories are not tracked.

\subsection{Congestion and Fleet Constraints}

Smart-locker service quality and population satisfy
\begin{align}
&
\sum_{c\in\mathcal C}\mu_c z_{ic}
-\sum_{u:(u,i)\in\mathcal A}\psi_{ui}
\ge
\frac{x_i}{\tau^H},
&&\forall i\in\mathcal H,
\label{app:locker_qos}
\\
&
\xi_i
\le M^3_i x_i,
&&\forall i\in\mathcal H,
\label{app:locker_population}
\\
&
\left\|
\begin{pmatrix}
2\displaystyle\sum_{c\in\mathcal C}\sqrt{\mu_c}\,z_{ic}
\\[1mm]
\xi_i
+\displaystyle\sum_{u:(u,i)\in\mathcal A}\psi_{ui}
-\displaystyle\sum_{c\in\mathcal C}\mu_c z_{ic}+1
\end{pmatrix}
\right\|_2
\le
\xi_i
-\displaystyle\sum_{u:(u,i)\in\mathcal A}\psi_{ui}
+\sum_{c\in\mathcal C}\mu_c z_{ic}+1,
&&\forall i\in\mathcal H.
\label{app:locker_queue_soc}
\end{align}
Constraint~\eqref{app:locker_qos} subtracts total inbound realized flow from the selected service rate to impose the mean handling-time requirement.
In constraint~\eqref{app:locker_queue_soc}, the first vector component represents the selected service rate in rotated-cone form, while the second combines expected population, inbound flow, and service rate.
The right-hand side completes the $M/M/1$ population epigraph after both rates have been substituted directly.
The virtual-state waiting requirement and fleet-size constraint are
\begin{align}
&
\sum_{r\in\mathcal R}
\frac{\rho_r}{1-\rho_r}\delta_{s_i r}
\le
\tau^S\psi_{s_i i},
&&\forall i\in\mathcal H,
\label{app:holding_qos}
\\
&
\sum_{\ell\in\widetilde{\mathcal L}}
\sum_{r\in\mathcal R}
\frac{\rho_r}{1-\rho_r}\delta_{\ell r}
+
\sum_{i\in\mathcal H}\xi_i
+
\sum_{(u,v)\in\mathcal A^{F}}
\sum_{k\in\widetilde{\mathcal E}_{uv}}
t_{uv}^{k}\varphi_{uv}^{k}
\le\Gamma.
\label{app:fleet_size}
\end{align}
In constraint~\eqref{app:holding_qos}, the left-hand side is the expected population of virtual state $s_i$, and the right-hand side is the maximum population allowed by its release rate and waiting-time limit.
In constraint~\eqref{app:fleet_size}, the three terms count drones in idle states, smart-locker handling queues, and physical transit, respectively.

\subsection{Mixed-Delivery Demand and Service Constraints}

The following constraints complete the mixed-delivery MLND formulation.
The platform may display at most one eligible destination locker, and any displayed locker must be one of the nearest eligible open lockers:
\begin{align}
&
\sum_{j\in\mathcal H_d}\alpha_{dj}
\le1,
&&\forall d\in\mathcal N,
\label{app:locker_destination_limit}
\\
&
\alpha_{dj}
\le x_j,
&&\forall d\in\mathcal N,\ j\in\mathcal H_d,
\label{app:locker_destination_open}
\\
&
\alpha_{dj}
\le1-x_{j'},
&&\forall d\in\mathcal N,\ j\in\mathcal H_d,\ j'\in\mathcal H_{dj}^{<}.
\label{app:nearest_locker}
\end{align}
Because at most one locker is displayed, the availability-adjusted MNL probability of choosing pickup at $j$ is exactly $\alpha_{dj}\eta_{odj}$, and the choice-stage home-delivery share is $1-\sum_{j\in\mathcal H_d}\alpha_{dj}\eta_{odj}$.
The auxiliary variable $\epsilon_{odj}$ represents successful market-$(o,d)$ flow ending at the displayed locker $j$ through the following exact linearization:
\begin{align}
&
\epsilon_{odj}
\le
\sum_{i:(i,j)\in\mathcal A^{HH}}
\sum_{r\in\mathcal R}
\rho_r\chi_{odijr},
&&\forall o,d\in\mathcal N,\ j\in\mathcal H_d,
\label{app:epsilon_flow}
\\
&
\epsilon_{odj}
\le M^4_{odj}\alpha_{dj},
&&\forall o,d\in\mathcal N,\ j\in\mathcal H_d,
\label{app:epsilon_display}
\\
&
\epsilon_{odj}
\ge
\sum_{i:(i,j)\in\mathcal A^{HH}}
\sum_{r\in\mathcal R}
\rho_r\chi_{odijr}
-M^4_{odj}(1-\alpha_{dj}),
&&\forall o,d\in\mathcal N,\ j\in\mathcal H_d.
\label{app:epsilon_lower}
\end{align}
The first inequality limits $\epsilon_{odj}$ by successful flow ending at $j$.
The second sets it to zero when $j$ is not displayed.
The third makes it equal that successful flow when $j$ is displayed.
The actual locker-pickup rate is $\eta_{odj}\epsilon_{odj}$, while the remaining successful flow uses last-mile courier delivery.
This representation assumes that the customer-choice outcome is independent of backend origin locker $i$, availability level $r$, and service realization conditional on $(o,d,j)$.

\subsection{Home-Delivery-Assistance Restriction}
\label{app:hd_model}

The home-delivery-assistance restriction inherits the infrastructure, OD assignment, availability, flow, routing, aggregate energy, congestion, holding, and fleet constraints of MLND.
It omits the customer-choice parameters, $\alpha_{dj}$, $\eta_{odj}$, and $\epsilon_{odj}$ because every realized drone order uses last-mile courier delivery.
Constraint~\eqref{app:od_demand_bound} allows the platform to offer any portion of each OD market, and constraints~\eqref{app:od_reliability_match} and~\eqref{app:realized_flow_definitions} retain the same forward service-realization logic.
All quantities not realized by the drone network are delivered directly from $o$ to $d$ by couriers.
Its complete objective is
\begin{align}
\min\quad
&
\sum_{u\in\mathcal L\cup\mathcal H\cup\mathcal E}f_u^{F}x_u
+
\sum_{i\in\mathcal H}\sum_{c\in\mathcal C}f_{ic}^{C}z_{ic}
+
\sum_{(u,v)\in\mathcal A^{Y}}f_{uv}^{A}y_{uv}
\nonumber\\
&+
\sum_{(u,v)\in\mathcal A^{F}}
\sum_{k\in\widetilde{\mathcal E}_{uv}}
f_{uvk}^{O}\varphi_{uv}^{k}
+f^{D}\Gamma
+
\sum_{o,d\in\mathcal N}f_{od}^{\mathrm{rider}}h_{od}
\nonumber\\
&+
\sum_{o,d\in\mathcal N}
\sum_{(i,j)\in\mathcal A^{HH}}
\sum_{r\in\mathcal R}
\rho_r
\left(
f_{oi}^{\mathrm{rider}}
+f_{jd}^{\mathrm{rider}}
-f_{od}^{\mathrm{rider}}
\right)
\chi_{odijr}.
\label{app:hd_objective}
\end{align}
The first line contains infrastructure costs.
The second line contains route, fleet, and all-courier baseline costs.
The final line replaces the direct courier cost only for orders realized by the drone network.

\clearpage
\subsection{Locker-Only Restriction}
\label{app:lo_model}

The locker-only restriction inherits constraints~\eqref{app:capacity_select}--\eqref{app:fleet_size}, including the aggregate energy-sufficiency constraints, and replaces only the mixed-delivery demand block with the constraints below.
It therefore admits only potential locker-pickup orders to the drone network.
Home delivery is the choice-stage outside alternative, with share $1-\sum_{j\in\mathcal H_d}\eta_{odj}\alpha_{dj}$ under the at-most-one displayed-locker rule.
The constraints follow.
\begin{align}
&
\sum_{j\in\mathcal H_d}\alpha_{dj}
\le1,
&&\forall d\in\mathcal N,
\label{app:lo_destination_limit}
\\
&
\alpha_{dj}\le x_j,
&&\forall d\in\mathcal N,\ j\in\mathcal H_d,
\label{app:lo_destination_open}
\\
&
\alpha_{dj}\le1-x_{j'},
&&\forall d\in\mathcal N,\ j\in\mathcal H_d,\ j'\in\mathcal H_{dj}^{<},
\label{app:lo_nearest_locker}
\\
&
\sum_{i:(i,j)\in\mathcal A^{HH}}
\sum_{r\in\mathcal R}
\chi_{odijr}
=h_{od}\eta_{odj}\alpha_{dj},
&&\forall o,d\in\mathcal N,\ j\in\mathcal H_d,
\label{app:lo_locker_allocation}
\\
&
\sum_{i:(i,j)\in\mathcal A^{HH}}
\sum_{r\in\mathcal R}
\chi_{odijr}
=0,
&&\forall o,d\in\mathcal N,\ j\in\mathcal H\setminus\mathcal H_d.
\label{app:lo_ineligible_locker}
\end{align}
Constraint~\eqref{app:lo_locker_allocation} assigns the complete choice-generated locker demand to backend origin lockers and availability levels before realization, while constraint~\eqref{app:lo_ineligible_locker} excludes corridors ending at a locker outside the customer's walking-distance set.
The reliability factor $\rho_r$ subsequently determines the realized locker-pickup quantity, and any unrealized quantity falls back to direct courier home delivery.
No $\epsilon_{odj}$ variable is needed because every successful drone order in this restriction completes locker pickup.
The complete objective is
\begin{align}
\min\quad
&
\sum_{u\in\mathcal L\cup\mathcal H\cup\mathcal E}f_u^{F}x_u
+
\sum_{i\in\mathcal H}\sum_{c\in\mathcal C}f_{ic}^{C}z_{ic}
+
\sum_{(u,v)\in\mathcal A^{Y}}f_{uv}^{A}y_{uv}
\nonumber\\
&+
\sum_{(u,v)\in\mathcal A^{F}}
\sum_{k\in\widetilde{\mathcal E}_{uv}}
f_{uvk}^{O}\varphi_{uv}^{k}
+f^{D}\Gamma
+
\sum_{o,d\in\mathcal N}f_{od}^{\mathrm{rider}}h_{od}
\nonumber\\
&+
\sum_{o,d\in\mathcal N}
\sum_{(i,j)\in\mathcal A^{HH}}
\sum_{r\in\mathcal R}
\rho_r
\left(
f_{oi}^{\mathrm{rider}}
+f_j^{\mathrm{locker}}
-f_{od}^{\mathrm{rider}}
\right)
\chi_{odijr}.
\label{app:lo_objective}
\end{align}
The first line contains infrastructure costs.
The second line contains route, fleet, and all-courier baseline costs.
The final line replaces the baseline cost of each realized locker order by its first-mile courier and locker-handling costs.

\endgroup

\end{document}
